% Modernised reading — fonds Alexandre Grothendieck, shelfmark 93
% (pages 1-235, the whole folder).
%
% This is an INTERPRETATION, not a transcription. It re-reads the pages in
% today's notation and vocabulary, states the hypotheses the manuscript leaves
% implicit, and drops the critical apparatus so that the mathematics reads
% straight through. Everything the brackets and underlines carried moves into a
% footnote. It is held to being mathematically correct as it stands; where the
% manuscript is loose or mistaken, the true statement is what appears here and
% the footnote says what the page has. In any dispute of substance the
% transcription governs: whoever cites Grothendieck must cite
% batch-01.fr.tex ... batch-12.fr.tex (pp. 1-20, 21-40, ..., 201-220, 221-235).
%
% Pass: Opus 5.5 (claude-opus-5-5), 2026-10-11 — first pass, unchecked against the pages by a human.
% The folder was read whole, all twelve batches; every batch is transcribed.
% It was made on Opus 5.5 and has not been measured against the reference
% reading of folder 115 (made on Opus 5). The literature named in the
% footnotes is cited from memory and was not checked against the sources.
% Models followed: the readings of folders 75, 82, 85 and 100 of the same
% group, cross-referenced where they meet this one.
% Revised: Opus 5.5 (claude-opus-5-5), 2026-10-18 — refined with the Lean
% checks of lean/Grothendieck/Folder93.lean (stations VI and IX): the
% Proposition of pp. 121-123 no longer normalises ad - bc = 1 (false, s = 2t
% over F_5; the determinant is a unit), and a footnote says Pic(A) = 0 is not
% used once alpha = 0 and that p = 2 needs no hypothesis on A; footnotes record
% the checks of G = 12 G_0, the cocycle without denominators, the p. 72 test
% values and the margin of p. 121 (level p^2 needed for p = 2, 3), the mathlib
% declarations and open PRs involved, and the two sources read by the finding
% 93-integral-schwarzian-level-p2 (Hoshi 2024, Wakabayashi 2021). No facsimile
% revision of the transcription postdates the reading.
% Revised: Opus 5.5 (claude-opus-5-5), 2026-10-18 — pp. 234-235: footnote on the
%   tier-2 Lean check of the derived morphism (Folder93.lean), for a torsor
%   under an abelian group: G-bar built without a point, [G -> G-bar] ~ Z,
%   (k, phi^0) a morphism of complexes, the class of u(p) independent of p;
%   the sheaf case and de Rham not checked. No literature (the finding is
%   unsearched). One « Ce que » turn removed.
%
% Reading is thin. The transcription carries about 12.6 \ill/\uncertain per
% page; on pp. 97-99, 109-112, 127-139, 188-197, 202-203, 205-217 whole
% sentences are guessed. The reading restates only what the transcription
% supports and says, section by section, where it is thin. Pages 138-139
% (invariant projective connections on torsors) and 215, 217 (struck) are
% summarised in a few lines only.
%
% Third-party material. Pages 2-24 are a copy of a typescript by P. Deligne,
% « Courbes elliptiques : formulaire, d'après J. Tate », annotated by a hand
% the transcription does not take for his; pages 27-29 are an offprint of
% M. Wanderley's C. R. note of 19 January 1981. Both are third-party work
% under RIGHTS.md: the transcription only notes them, and this reading says
% no more than its notes (titles, section headings, where the annotations
% are). Their mathematics is not restated.
%
% What the folder is. Not one argument but a box of separate dossiers, each
% under its own cover, which share one theme: differential structures on
% families — connections and their integrability, in characteristic 0 and p.
% Spine (stations, with pages):
%   I     pp. 2-24     Deligne's typescript (third party, described only)
%   II    pp. 26-32    axioms for a theory of orientation (after Wanderley)
%   III   p. 35        an isolated Schwarzian computation
%   IV    pp. 36-53    his typescript « Programme de la théorie de Dieudonné »
%                      (crystalline Dieudonné functor, deformation theorems,
%                      admissible F-V-crystals; « rabiots », Problems 1-6,
%                      Tate's theorem, the « foncteur mystérieux »)
%   V     pp. 54-67    chi and chi_! (Euler characteristics with and without
%                      compact supports; locally conic spaces); theme list
%   VI    pp. 71-83    differential invariants of étale maps of curves: the
%                      Schwarzian G as the unique invariant; other groups
%                      (translations, homotheties, affine) ; p. 82 Postnikov
%   VII   pp. 85-100   P^1-bundles, sl_2-type Lie algebras, transverse
%                      connections (sl_2-opers), horizontal sections
%   VIII  pp. 101-114  c_{D}, change of coordinate, horizontality equation
%                      2 lambda' + lambda^2 = mu, the dictionary (a)-(e)
%   IX    pp. 115-123  fair copy: G(t,s) cocycle, torsor of projective
%                      connections, relative case, the factor 12, char p
%   X     pp. 125-139  uniformisation and the canonical projective
%                      connection; regular singularities; counts M_{g,nu}
%   XI    pp. 141-159  connections and stratifications on affine torsors;
%                      Picard-Fuchs jets and operators, « section de Manin »
%   XII   pp. 161-166  filtered de Rham complex and Gauss-Manin (Katz-Oda)
%   XIII  pp. 167-185  the Cartier operation (p-bases, log forms, trace,
%                      residues, duality, Verschiebung, curves: pairing)
%   XIV   pp. 187-197  problems: Cartier descent, stratified modules and
%                      Pi_1, the p-curvature question (« Fin 1ère partie »)
%   XV    pp. 198-203  iterated Cartier operation, Frobenius eigenvalues
%   XVI   pp. 204-212  relative Picard functors with connection P^c ... P^cin
%   XVII  pp. 213-217  c^1(L), integrable connections, p-curvature
%   XVIII pp. 218-220  plan for the canonical Gauss-Manin connection
%   XIX   pp. 221-229  integrals of closed relative forms; Picard-Fuchs
%   XX    pp. 230-232  differential operators with constant coefficients
%   XXI   pp. 233-235  canonical isomorphism of de Rham cohomologies of two
%                      liftings (homotopy formula), runs past the folder's end
%
% Notation fixed for the whole document. omega = Omega^1_{X/S} for a curve
% (his underlined omega); omega^{n} its n-th tensor power. A projective
% connection is a section of the torsor Conn_{X/S} under omega^{2};
% G(t,s) = c(s) - c(t) the Schwarzian difference (his G); G_0 = G/12. In
% VII-VIII, E is the rank-2 extension 0 -> Omega -> E -> O -> 0 (his \mathcal{L}
% on pp. 86-100, his E on p. 132), g the Lie algebra (form of pgl_2),
% b subset g the Borel subalgebra (also \mathcal{L} on the pages; b ~ E as
% extensions, p. 86). X' = X^{(p/S)}, F_{X/S} relative Frobenius, C-tilde the
% Cartier operation. P^infty_{S/T}(E) the module of infinite jets.
%
% Substantive departures from the pages, each footnoted where it occurs:
% — p. 30: the « rives » sheaf is a double cover of Y, not of X; p. 32, 4°)
%   « omega in Gamma(X) » read omega in Omega(X).
% — p. 35: G(t, exp f) = G(t,f) + (df)^2, coefficient 1 (the page's 3 is
%   wrong; its first coefficient is illegible).
% — p. 43, 48: Fil^1 = Ker F = Im V up to the conventions of the
%   contravariant theory; stated as the page has it.
% — p. 57: chi_!(R^n) for chi(R^n); p. 61 E_{n-1} read C_{n-1}; p. 61
%   marginal: compact 4-dim locally conic spaces need not be triangulable.
% — p. 72: G_0(t,t^2) = dt^2/(4t^2), not dt^2/(2t).
% — p. 81: sign convention of beta_n fixed as on p. 79.
% — pp. 96, 99, 102: the factor 2 of the horizontality equation (he flags it
%   himself; corrected p. 105): 2 lambda' + lambda^2 = mu.
% — p. 113: G_tau(s) = (G_t(s) - G_t(tau))/tau'^2, not G_t(s) tau'^2.
% — p. 114: P^1 minus three points is not a doubtful case (p. 137 settles it).
% — p. 125: the stabiliser of a disc is PSL_2(R), the identity component of
%   PGL_2(R); the correspondence real forms <-> discs holds for split forms.
% — p. 135: « régulière à l'infini » read with fixed (cuspidal) leading term,
%   which is what makes the group H^0(omega^2(D)); general regular singular
%   projective connections form a torsor under omega^2(2D).
% — p. 182: C(x^i d x^j) nonzero iff i + j = p; p. 183 the sign of the
%   antidiagonal determinant is (-1)^{(p-1)(p-2)/2}, not (-1)^{p-2}.
% — p. 203: « equal to the number of eigenvalues congruent to 1 » holds as an
%   inequality in general; equality needs semisimplicity.
% — p. 208: the p-curvature lives in End(E) (x) F^*Omega^1_{X'}, which is
%   Omega^{1 (x) p} only in relative dimension 1.
% — p. 233-234: the homotopy written additively, u(pg) - u(p) = dk + kd.
%
% Modern names given and footnoted as ours: Grothendieck-Messing deformation
% theory, crystalline Dieudonné theory, universal vector extension,
% « foncteur mystérieux » and Fontaine's theory, Tate's theorem (de Jong in
% char p), Euler characteristic with compact supports, Schwarzian derivative,
% projective structure, sl_2-oper, Riccati equation, accessory parameters,
% Takhtajan-Zograf, Gauss-Manin connection, Katz-Oda, Picard-Fuchs equation,
% Manin map, Cartier isomorphism, Tate's residue formula, Hasse-Witt and
% Cartier-Manin, conjugate filtration, canonical theta characteristic in
% char 2, Cartier descent, p-curvature, Grothendieck-Katz conjecture,
% stratified bundles (Gieseker, dos Santos, Esnault-Mehta), Mazur-Messing.
%
% What the folder announces and never establishes: Problems 1-6 of the
% typescript; the uniqueness of chi, chi_! on all of E_infty (p. 65);
% formula (15) of p. 227 (admitted); the independence of the choice of
% p-basis on p. 184 (sketch); the Verschiebung statement of p. 180
% (admitted); the questions of pp. 193-197; the last construction (p. 235)
% runs past the folder's end.
\documentclass[11pt,a4paper]{article}
\input{../preamble/grothendieck}

\folder{93}
\pages{1}{235}
\dating{[1972]-1981}
\watermark{Édition de démonstration\\interprétation personnelle de l'œuvre}
\foldertitle{Connexions de Gauss-Manin et équations de Picard-Fuchs. Opération de Cartier : copies de tapuscrit annoté (s.d.), tiré à part (1981), notes manuscrites (s.d.) — lecture modernisée du dossier entier}

\begin{document}

\section*{Résumé}

\begin{resume}
Une équation différentielle dont les coefficients dépendent d'un paramètre a
des solutions qui dépendent, elles aussi, de ce paramètre ; et la manière dont
elles en dépendent obéit à son tour à une équation différentielle. C'est le
phénomène que ce dossier poursuit sous plusieurs formes. Le plus ancien
exemple est celui des intégrales elliptiques : leurs périodes, vues comme
fonctions du module de la courbe, satisfont une équation linéaire du second
ordre, l'\emph{équation de Picard--Fuchs}. L'objet géométrique qui la porte
s'appelle aujourd'hui la \emph{connexion de Gauss--Manin} : une règle qui dit
comment « transporter » une classe de cohomologie d'une fibre d'une famille à
la fibre voisine.

Le dossier n'est pas un texte suivi mais une boîte de chemises, chacune avec
sa couverture et son titre. Ce qui les relie est la notion de
\emph{connexion} — une façon de dériver dans une direction donnée — et la
question de savoir quand une connexion est \emph{intégrable}, c'est-à-dire
quand elle provient d'un système de solutions. Le dossier pose cette question
en caractéristique nulle, où l'analyse complexe guide l'intuition, et en
caractéristique $p$, où tout change : la dérivée $p$-ième d'une fonction est
nulle, la fonction $x^p$ a une dérivée nulle sans être constante, et un
outil propre à cette situation, l'\emph{opération de Cartier}, prend la place
de l'intégration.

On y trouve : un programme dactylographié de la
théorie de Dieudonné « cristalline », qui décrit les groupes $p$-divisibles
par des modules munis d'un Frobenius et d'une connexion ; deux chemises plus
courtes, l'une sur les axiomes d'une théorie de l'orientation des variétés,
l'autre sur la caractéristique d'Euler à supports compacts ; une longue étude
de la \emph{dérivée schwarzienne} et des \emph{connexions projectives} sur une
courbe ; des notes sur la connexion de Gauss--Manin et les équations de
Picard--Fuchs ; un exposé de l'opération de Cartier ; et des questions. Une
copie d'un formulaire sur les courbes elliptiques de P.~Deligne, d'après
J.~Tate, ouvre le dossier ; c'est le travail d'un autre, on ne fait que la
signaler.

La dérivée schwarzienne est la pièce la plus lisible pour qui découvre le
sujet. Si l'on change de coordonnée sur une droite par une homographie
$t \mapsto (at+b)/(ct+d)$, aucune des dérivées $s'$, $s''$, $s'''$ d'une fonction
$s$ n'est conservée, mais une certaine combinaison d'entre elles l'est. Les
pages montrent que cette combinaison est, à un facteur près, la \emph{seule}
qui soit invariante, et qu'elle mesure l'écart entre deux façons de mettre
sur une courbe une « structure projective » : un atlas dont les changements
de cartes sont des homographies. Le théorème d'uniformisation dit qu'une
telle structure existe canoniquement sur toute surface de Riemann ; les
feuillets demandent quand elle est algébrique, et comptent, avec les
dimensions des espaces de modules de courbes, combien il y a de structures
« régulières » possibles.

Ce qui frappe, en caractéristique $p$, est la question posée aux pages 193
et 194. Une connexion sur une famille de variétés définie sur les entiers se
réduit modulo chaque nombre premier ; si elle provient de solutions
algébriques, sa « $p$-courbure » est nulle pour presque tout $p$. La page
demande si la réciproque est vraie. C'est la question qu'on appelle
aujourd'hui la \emph{conjecture de Grothendieck--Katz sur la $p$-courbure} ;
elle est toujours ouverte.

On notera un geste qui revient. Devant l'expression $-2\,s'''/s' +
3\,(s''/s')^2$, Grothendieck ne s'arrête pas à sa formule : il demande ce
qu'elle devient sur les entiers, trouve que ses coefficients $2$ et $3$
l'empêchent d'y vivre, et la divise par $12$ ; la formule ainsi normalisée
a un sens en toute caractéristique, et l'invariant n'est plus déterminé
qu'au signe près. L'objet canonique n'est pas la formule qu'on connaît,
mais celle qui survit au changement de base.

\keywords{Gauss-Manin connection, Picard-Fuchs equation, Katz-Oda filtration, Cartier operator, Cartier isomorphism, logarithmic differential, p-curvature, Cartier descent, Grothendieck-Katz p-curvature conjecture, stratification, stratified bundle, crystalline Dieudonné theory, Grothendieck-Messing deformation theory, universal vector extension, mysterious functor, Schwarzian derivative, projective connection, projective structure, sl2-oper, Riccati equation, uniformization theorem, accessory parameters, Euler characteristic with compact supports, locally conic space, orientation of topological manifolds, Manin map, Hasse-Witt matrix, conjugate filtration, theta characteristic, Picard functor with connection}
\end{resume}

\section*{Le fil du dossier, et les conventions}

\subsection*{Les stations}

L'inventaire date le dossier « [1972]-1981 » et le range dans le groupe
« Géométrie et topologie combinatoire » (dossiers 69 à 102), dont le sujet
n'est, ici encore, pas le sien.\footnote{Comme celui du dossier 100, qui porte
sur la bonne réduction des schémas abéliens et rencontre celui-ci à la
station IV. Le millésime 1981 est celui du tiré à part de M.~Wanderley
(pages 27 à 29) ; aucune page de sa main n'est datée, et on ne tire rien des
dates des versos.} Le dossier est fait de chemises successives ; chacune a sa
couverture, souvent de sa main, et l'ordre des feuillets est celui des
archivistes. On le suit.

\begin{enumerate}
  \item[I.] Pages 2 à 24 : copie annotée d'un tapuscrit de P.~Deligne,
  « Courbes elliptiques : formulaire, d'après J.~Tate » (texte d'un tiers,
  signalé seulement).
  \item[II.] Pages 26 à 32 : « Orientation » — un tiré à part de M.~Wanderley
  (1981), puis ses axiomes pour une théorie de l'orientation.
  \item[III.] Page 35 : un calcul de dérivée schwarzienne, isolé.
  \item[IV.] Pages 36 à 53 : « Programme de la théorie (Cristaux) » — son
  tapuscrit sur la théorie de Dieudonné, corrigé de sa main, et ses
  « rabiots ».
  \item[V.] Pages 54 à 67 : « $\chi$ et $\chi_!$ » — caractéristiques d'Euler
  avec et sans supports compacts.
  \item[VI.] Pages 71 à 83 : invariants différentiels des applications
  étales d'une courbe dans la droite ; la dérivée schwarzienne comme seul
  invariant projectif ; les autres groupes.
  \item[VII.] Pages 85 à 100 : fibrés en droites projectives, algèbres de Lie
  de type $\mathfrak{sl}_2$, connexions transverses à une section.
  \item[VIII.] Pages 101 à 114 : changement de coordonnée, équation
  d'horizontalité, dictionnaire des structures équivalentes.
  \item[IX.] Pages 115 à 123 : mise au net — le cocycle schwarzien, le
  torseur des connexions projectives, le cas relatif, le facteur $12$, la
  caractéristique $p$.
  \item[X.] Pages 125 à 139 : uniformisation et connexion projective
  canonique ; singularités régulières ; connexions invariantes.
  \item[XI.] Pages 141 à 159 : connexions et stratifications sur les torseurs
  affines ; équations de Picard--Fuchs (couverture page 160).
  \item[XII.] Pages 161 à 166 : filtration du complexe de De~Rham et connexion
  de Gauss--Manin, d'après Katz--Oda.
  \item[XIII.] Pages 167 à 186 : l'opération de Cartier (couverture page 167).
  \item[XIV.] Pages 187 à 197 : « Connexions et stratifications intégrables :
  (suite, problèmes divers) » ; la page 197 dit « Fin 1ère partie ».
  \item[XV.] Pages 198 à 203 : itérés de l'opération de Cartier et valeurs
  propres de Frobenius (« Fin de parenthèse »).
  \item[XVI.] Pages 204 à 212 : « Variantes concernant des foncteurs de
  Picard relatifs ».
  \item[XVII.] Pages 213 à 217 : « $c^1(L)$ style connexions et opération de
  Cartier ».
  \item[XVIII.] Pages 218 à 220 : « Connexion canonique de Gauss--Manin »,
  un plan.
  \item[XIX.] Pages 221 à 229 : intégrales de formes fermées relatives et
  équations de Picard--Fuchs.
  \item[XX.] Pages 230 à 232 : opérateurs différentiels à coefficients
  constants.
  \item[XXI.] Pages 233 à 235 : construction de la connexion canonique en
  cohomologie de De~Rham ; la construction se poursuit au-delà du dossier.
\end{enumerate}

Trois fils traversent ces chemises. Le premier est celui des
\emph{connexions projectives} (stations III, VI à X) : il commence par un
calcul isolé et finit par une théorie sur une base quelconque. Le deuxième
est celui de la \emph{connexion de Gauss--Manin et des équations de
Picard--Fuchs} (stations XI, XII, XVIII, XIX, XXI), qui en donne trois
constructions : par la filtration de Katz--Oda, par des torseurs, par
l'intégration sur des chemins. Le troisième est celui de la
\emph{caractéristique $p$} (stations IV, XIII à XVII) : Frobenius, Cartier,
$p$-courbure. Les feuillets ne se renvoient presque jamais les uns aux autres ;
les renvois qu'on fait ici entre stations sont notre lecture.

\subsection*{Conventions, valables pour tout le dossier}

Pour une courbe lisse $X$ sur $S$, on note
$\omega = \Omega^1_{X/S}$ (il souligne : $\underline{\omega}$) et $\omega^n$ sa
puissance tensorielle $n$-ième. Une \emph{uniformisante} est une fonction
$t$ telle que $dt$ soit une base de $\omega$, c'est-à-dire une application
étale vers la droite affine ; pour deux uniformisantes $t$, $s$ on écrit
$s' = ds/dt$, $s'' = ds'/dt$, etc. On note
\[
  G(t,s) = \Bigl(-2\,\frac{s'''}{s'} + 3\Bigl(\frac{s''}{s'}\Bigr)^2\Bigr)\,dt^2
  = -2\,\{s,t\}\,dt^2,
\]
où $\{s,t\} = s'''/s' - \tfrac32 (s''/s')^2$ est la dérivée schwarzienne
classique ;\footnote{La notation $\{s,t\}$ et le nom de Schwarz ne figurent
pas sur ces pages : il écrit $G$, $G_t(s)$ ou $C(t,s)$ (page 35), et ne nomme
pas l'expression. Le facteur $-2$ qui la sépare de la normalisation
classique est le sien ; on le garde partout, car toutes les formules du
dossier sont écrites avec lui.} $G_0 = G/12$ est la forme normalisée de la
page 120. Une \emph{connexion projective} est une section du torseur
$\mathrm{Conn}_{X/S}$ sous $\omega^2$ (pages 117 et 119), et
$G(t,s) = c(s) - c(t)$.

Aux stations VII et VIII, $E$ désigne le module localement libre de rang $2$
extension de $\mathcal{O}_X$ par un module inversible $\Omega$,
$0 \to \Omega \to E \to \mathcal{O}_X \to 0$ ; $\mathfrak{g}$ l'algèbre de Lie
associée (une forme de $\mathfrak{pgl}_2$), $\mathfrak{b} \subset \mathfrak{g}$
la sous-algèbre de Borel, $P = \mathbb{P}(E)$ le fibré en droites projectives
et $\sigma$ sa section.\footnote{Les pages 86 à 100 notent $\mathcal{L}$ à la
fois le module de rang $2$ et la sous-algèbre de Borel ; la page 86 établit
que les deux sont isomorphes comme extensions de $\mathcal{O}_X$ par
$\Omega$, ce qui explique l'usage. On les distingue ici. La page 132 écrit
$E$ pour le même module ; on reprend sa lettre.} En caractéristique $p$, on
note $X' = X^{(p/S)}$ le tordu de Frobenius de $X$ sur $S$,
$F_{X/S} : X \to X'$ le Frobenius relatif et $\widetilde{C}$ l'opération de
Cartier. $P^n_{S/T}(E)$ est le module des jets d'ordre $n$ de $E$, et
$P^\infty_{S/T}(E)$ leur limite projective.

Deux homonymies à ne pas confondre. La lettre $P$ désigne un fibré en
droites projectives aux stations VII à X et un torseur aux stations XI et
XVI ; on rappelle chaque fois lequel. Et la lettre $\mathcal{F}$ désigne aux
pages 223 à 227 le quotient $f_*(Z\Omega^1_{X/S})/d f_*(\mathcal{O}_X)$, et
page 235 un faisceau quelconque.

\pagerange{2}{24}
\section*{I. Un formulaire sur les courbes elliptiques, d'un autre (pages 2 à 24)}

Les pages 2 à 24 sont la copie d'un tapuscrit qui n'est pas de Grothendieck :
« Courbes elliptiques : formulaire, d'après J.~Tate, mis au goût du jour par
P.~Deligne », sans date. C'est le travail d'un tiers, que la transcription ne
reproduit pas ; on n'en dit que ce que disent ses notes. Les titres de
paragraphes relevés sont : « Formules générales » (§~1), « Cas où $2$ et $3$
sont inversibles » (§~2), « Cas où $2$ et $c_4$ sont inversibles » (§~3), « Cas
où $2$ est nilpotent et $c_4$ inversible » (§~4), « Cas où $c_4 = 0$, $p = 2$ ou
$p = 3$ » (§~5), « Applications » (§~6), « Anneau des formes modulaires
entières » (§~7).\footnote{La note de la page 2 dit le tapuscrit paginé de 1
à 16 ; celle de la page 21 lui donne des pages 17 à 19, et la page 24 est
dite « p.~19 de la dactylographie ». C'est la seconde indication qui
s'accorde avec le nombre des feuillets.}

La copie porte des annotations à l'encre noire, verte et bleue, d'une main
qui écrit en tête de la page 2 « G? = correction de Grothendieck » et renvoie
à une version plus complète dans les \emph{Lecture Notes} 476 (Anvers IV). Cette
main parle de lui à la troisième personne ; la transcription ne la prend pas
pour la sienne, et l'attribution à Grothendieck des corrections marquées
« G? » est la sienne, interrogative. Deux versos portent des notes de cette
même main (pages 17 et 20), une page manuscrite de « Formules utiles » est
intercalée (page 3) ; on ne les restitue pas davantage.

\pagerange{26}{32}
\section*{II. Axiomes d'une théorie de l'orientation (pages 26 à 32)}

Un feuillet blanc porte le seul mot « Orientation ». Suivent un tiré à part,
M.~Wanderley, « Théorie axiomatique de l'orientation des variétés
topologiques à bord », note présentée par H.~Cartan, \emph{C.~R. Acad. Sc.
Paris}, t.~292, série~I, p.~213--215, 19 janvier 1981 (pages 27 à 29, sans
annotation de sa main, non transcrites), puis deux pages de sa main qui
posent des axiomes du même genre. On ne peut pas dire, faute du texte de la
note, en quoi les siens s'en écartent.

\textbf{Le cadre.} Soit $\overline{\mathcal{V}}$ la catégorie des variétés
topologiques à bord, avec pour morphismes les immersions ouvertes, et
$\mathcal{V} \subset \overline{\mathcal{V}}$ la sous-catégorie des variétés sans
bord ; $\partial X$ est le bord de $X$, $\operatorname{Int} X = X - \partial X$
son intérieur. Une \emph{théorie de l'orientation} est la donnée :
\begin{enumerate}
  \item[(i)] d'un foncteur contravariant
  $\Omega : \mathcal{V}^{\circ} \to (\mathrm{Ens})$, $\Omega(X)$ étant
  l'ensemble des orientations de $X$ ;
  \item[(ii)] d'un élément $\omega^+ \in \Omega(e)$, où $e$ est la variété
  ponctuelle ;
  \item[(iii)] d'une transformation naturelle, entre foncteurs sur
  $\overline{\mathcal{V}}^{\circ}$, de $X \mapsto \Omega(\operatorname{Int} X)$
  vers $X \mapsto \Omega(\partial X)$ — la règle qui oriente le bord d'une
  variété orientée ;
\end{enumerate}
soumise aux conditions suivantes.
\begin{enumerate}
  \item[1°)] $\Omega$ est un faisceau : sa restriction aux ouverts de chaque
  $X$ est un faisceau $\Omega_X$, le \emph{faisceau d'orientation} de $X$.
  \item[2°)] $\Omega(e)$ a deux éléments. On note $\omega^-$ l'autre, et l'on
  identifie $\Omega(e)$ à $\{+1,-1\}$.
\end{enumerate}
Les variétés munies d'une section de $\Omega$ forment la catégorie des
variétés $\Omega$-orientées. Si $X$ est orientée et $Y \subset X$ une
sous-variété de codimension $1$, le faisceau $\rho_{Y/X}$ des \emph{rives} de
$Y$ dans $X$ est un revêtement double de $Y$ ;\footnote{La page écrit « rev.
étale d'ordre 2 de $X$ », lecture incertaine du dernier mot ; c'est un
revêtement de $Y$ (au voisinage d'un point de $Y$, les deux côtés), comme
le veut la formule qui suit. On suppose, ce que la page ne dit pas, $Y$
localement plate, pour que les rives aient un sens.} la donnée (iii),
appliquée aux demi-voisinages de $Y$, fournit un morphisme
\[
  \Omega_X|_Y \times \rho_{Y/X} \longrightarrow \Omega_Y ,
  \qquad\text{ou}\qquad
  \Omega_X|_Y \longrightarrow \underline{\mathrm{Hom}}(\rho_{Y/X}, \Omega_Y).
\]
Les conditions se poursuivent :
\begin{enumerate}
  \item[3°)] (\emph{hypothèse d'induction}) ce dernier morphisme est
  injectif : une orientation de $X$ en un point de $Y$ est connue quand on
  sait quelle orientation de $Y$ chaque rive détermine ;
  \item[4°)] (\emph{non-trivialité}) pour toute orientation $\omega$ de $X$,
  le morphisme $\rho_{Y/X} \to \Omega_Y$ correspondant est injectif : deux
  rives différentes induisent des orientations différentes ;\footnote{La page
  écrit « pour tte $\omega \in \Gamma(X)$ » ; on lit $\omega \in \Omega(X)$.}
  \item[5°)] pour toute $X$ sans bord, $\Omega_X \to X$ est surjectif — il
  suffit que $\Omega(\mathbb{R}^n)$ soit non vide pour $n \geq 1$ ;
  \item[6°)] (\emph{prolongement des identités}) $\Omega_X$, vu comme espace
  étalé sur $X$, est séparé : une orientation est connue en un point quand
  on la connaît en des points arbitrairement voisins.
\end{enumerate}
Il ajoute, d'une autre encre, qu'il suffirait d'exiger 6°) pour les
$\mathbb{R}^n$, et qu'« il sera vrai a posteriori » qu'une orientation de
$\mathbb{R}^n$, ou de toute $X$ connexe, est connue quand on la connaît en un
seul point. La page s'arrête là ; aucun théorème n'est énoncé. Le but des
axiomes est visible : 1°), 5°) et 6°) font de $\Omega_X$ un revêtement
de $X$, 2°) à 4°) forcent par récurrence sur la dimension ses fibres à avoir
deux éléments, et l'on retrouve le revêtement d'orientation, torseur sous
$\{\pm 1\}$.\footnote{Cette conclusion est la nôtre ; les feuillets ne la
démontrent pas. Les orientations comme torseurs sous $\{\pm1\}$ sont le sujet
des pages 24 à 27 du dossier 75 (cours C4 de 1977-1978), dans un cadre
combinatoire.}

\pagerange{35}{35}
\section*{III. Un calcul isolé (page 35)}

Une feuille de calculs sans texte suivi calcule l'effet d'une exponentielle
sur l'expression $-2\,g'''/g' + 3\,(g''/g')^2$. Pour $g = \exp f$ on a
$g' = (\exp f)\,f'$, $g'' = (\exp f)(f'^2 + f'')$,
$g''' = (\exp f)(f'^3 + 3f'f'' + f''')$, d'où
\[
  -2\,\frac{g'''}{g'} + 3\Bigl(\frac{g''}{g'}\Bigr)^2
  = -2\,\frac{f'''}{f'} + 3\Bigl(\frac{f''}{f'}\Bigr)^2 + f'^2 ,
\]
soit, dans les notations des conventions,
\[
  G(t, \exp f) = G(t, f) + (df)^2 .
\]
\footnote{Le coefficient de $f'^2$ est sur la page un petit signe superposé
que la transcription ne lit pas ; les deux lignes suivantes écrivent
$3(df)^2$, et l'encadré « $c_{\exp f} = c_f + \ldots\, df^2$ » laisse le
coefficient illisible. Le calcul donne $1$ ($3 - 2$), comme la transcription
le note.} C'est la formule $c_{D + \omega} = c_D + 2D\omega + \omega^2$ des
pages 81 et 104 dans un cas particulier (station VIII). Au verso, écrit la
feuille retournée, le même objet est noté $C(z,f)$, avec deux annotations :
« invariance par $\pi$ », reliée à $f$, et « $df/dz$ forme inv. par $\pi$ » ;
les seconds membres des dernières égalités, sur l'action d'une transformation
$\gamma$, ne se lisent pas. On n'en tire rien de plus.

\pagerange{36}{47}
\section*{IV. Programme de la théorie de Dieudonné (pages 36 à 53)}

Le titre « V Programme de la Théorie (Cristaux) » est écrit de sa main sur un
verso de listing (page 36) ; suit un tapuscrit de sa main, corrigé à la
main, marqué en marge « Nouveau chapitre » et « $\mathbb{D}^*$ ». Le chapitre
appartient à un exposé plus long, auquel il renvoie (« cf.\ plus haut »,
« dans la suite du séminaire ») sans que les renvois soient remplis.

\subsection*{1) Le foncteur de Dieudonné cristallin (pages 37 à 40)}

Soit $S$ un schéma sur lequel $p$ est localement nilpotent, et
$S_0 = V(p) \subset S$. On note $\mathrm{BT}(S)$ la catégorie des groupes de
Barsotti--Tate (groupes $p$-divisibles) sur $S$ et $\mathrm{Crisloclib}(S)$
celle des cristaux en modules localement libres sur le site cristallin de
$S$. Le programme est de définir un foncteur contravariant
\[
  \mathbb{D}^* : \mathrm{BT}(S)^{\circ} \longrightarrow \mathrm{Crisloclib}(S),
\]
compatible aux images inverses. Comme la restriction
$\mathrm{Crisloclib}(S) \to \mathrm{Crisloclib}(S_0)$ est une équivalence (l'idéal
$p\mathcal{O}_S$ est muni de puissances divisées), $\mathbb{D}^*(G)$ ne dépend
que de $G_0 = G \times_S S_0$, et l'on est ramené à une base de
caractéristique $p$.

Sur $S_0$ de caractéristique $p$, le Frobenius et le Verschiebung de $G_0$ et
la compatibilité aux images inverses munissent $M = \mathbb{D}^*(G_0)$ de
morphismes
\[
  F_M : M^{(p)} = F_{S_0}^*(M) \longrightarrow M, \qquad
  V_M : M \longrightarrow M^{(p)}, \qquad
  V_M F_M = p, \quad F_M V_M = p .
\]
Un tel triple est un \emph{F-V-cristal}, ou \emph{cristal de Dieudonné}, et le
foncteur $G_0 \mapsto (M, F_M, V_M)$ est le \emph{foncteur de Dieudonné}. Sa
formation commute à tout changement de base, et sur le spectre d'un corps
parfait il coïncide avec le foncteur de Dieudonné classique, qui est alors
une équivalence ; le cristal $\mathbb{D}^*(G_0)$ « saisit » donc les fibres
de $G_0$ en tous les points parfaits.

\emph{Problème 1.} Sur une base $S_0$ de caractéristique $p$, le foncteur de
Dieudonné est-il pleinement fidèle, et son image essentielle est-elle formée
des F-V-cristaux « admissibles » (définis plus bas) ?

Il le croit lié à une théorie de Dieudonné pour les schémas en groupes finis
localement libres d'ordre une puissance de $p$, qui demanderait des
« structures supplémentaires » — « très probablement » une filtration sur un
objet convenable d'une catégorie dérivée —, sauf peut-être sur une base
parfaite.

\emph{Problème 2.} Développer une telle théorie ; sur $S_0$ parfait, établir
une anti-équivalence entre les groupes finis localement libres sur $S_0$
plats sur $\Lambda_n = \mathbb{Z}/p^n$ et les $W_n(\mathcal{O}_{S_0})$-modules
localement libres munis de $F$ et $V$.\footnote{Ces deux problèmes ont été au
centre de la théorie de Dieudonné cristalline développée ensuite (Messing,
Mazur et Messing, Berthelot, Breen ; A.~J.~de~Jong, 1995, pour la pleine
fidélité sur certaines bases). Nous ne précisons pas ici l'état de chacun ; le nom
« Barsotti--Tate » est le sien (congrès de Nice, 1970), et la page l'emploie.}

\textbf{La filtration de Hodge.} Sur $S$ quelconque ($p$ localement
nilpotent), la valeur $\mathbb{D}^*(G)_S$ du cristal sur $S$ porte une
filtration par un sous-module localement facteur direct
$\mathrm{Fil}^1 = \omega_G$, d'où une suite exacte fonctorielle et compatible
aux changements de base
\[
  0 \longrightarrow \omega_G \longrightarrow \mathbb{D}^*(G)_S
  \longrightarrow \mathfrak{t}_{G^*} \longrightarrow 0,
\]
où $\omega_G$ est le module des différentielles invariantes du groupe formel
associé à $G$ et $\mathfrak{t}_{G^*} = \omega_{G^*}^{\vee}$ l'algèbre de Lie du
dual de Cartier--Serre $G^*$. On a ainsi un foncteur vers les « F-V-cristaux
filtrés ». C'est, dit la remarque de la page 41, l'analogue d'une filtration
de Hodge sur un $H^1$ de De~Rham.

\pagerange{41}{43}
\subsection*{Le théorème de déformation (pages 41 et 42)}

Soit $S \subset S'$ un épaississement muni de puissances divisées
\emph{nilpotentes}. Pour tout relèvement $G'$ de $G$ à $S'$, on a
$\mathbb{D}^*(G')_{S'} = \mathbb{D}^*(G)_{S'}$ (valeur du cristal sur l'objet
$S'$ du site), et la filtration de Hodge de $G'$ relève celle de $G$.

\textbf{Théorème de déformation.} Le foncteur qui à $G' \in \mathrm{BT}(S')$
associe le couple formé de $G = G'|_S$ et de la filtration
$\mathrm{Fil}^1 \subset \mathbb{D}^*(G)_{S'}$ relevant $\omega_G$ est une
équivalence de catégories.\footnote{L'hypothèse de nilpotence des puissances
divisées est un ajout manuscrit, dont la fin est coupée par le bord ; elle est
nécessaire, et c'est sous elle que le théorème a été démontré par
W.~Messing (\emph{The crystals associated to Barsotti--Tate groups},
LNM 264, 1972) : on l'appelle aujourd'hui théorème de Grothendieck--Messing.
Une note marginale, sur le cas des groupes infinitésimaux, ne se lit pas.}

Comme $(S_0, S)$ est lui-même un tel épaississement, une réponse positive au
problème 1 décrirait $\mathrm{BT}(S)$ entièrement par des cristaux de
Dieudonné filtrés ; et le cas général $(S, S')$ se ramène formellement au
cas $(S_0, S')$. \emph{Problème 3} (ajouté à la main) : une variante pour
les groupes finis localement libres — « j'ai un candidat ».

\textbf{F-V-cristaux admissibles.} En caractéristique $p$ la filtration de
Hodge est déterminée par la structure de cristal de Dieudonné. Soit
$\phi : S_{0,\mathrm{cris}/\mathbb{F}_p} \to S_0$ le morphisme de topos annelés,
$M_0$ la restriction de $M$ au site cristallin relatif à $\mathbb{F}_p$. Les
relations $FV = VF = p$ rendent nuls les composés de
$M_0^{(p)} \xrightarrow{F} M_0 \xrightarrow{V} M_0^{(p)} \xrightarrow{F} M_0$, et
la suite est même exacte ; et l'on a
\[
  \phi^*(\mathrm{Fil}^1 M_{S_0}) = \operatorname{Ker} F_{M_0}
  = \operatorname{Im} V_{M_0} \subset M_0^{(p)} .
\]
Une propriété de fidélité de $\phi$ fait que cette relation détermine
$\mathrm{Fil}^1$. On dit qu'un F-V-cristal est \emph{admissible} s'il existe un
$\mathrm{Fil}^1$ localement facteur direct vérifiant cette relation (il est
alors unique : c'est la \emph{filtration canonique}) ; le cristal d'un groupe
de Barsotti--Tate est admissible, et sa filtration de Hodge est la
filtration canonique.\footnote{Ces énoncés sont annoncés (« on prouvera »,
« on montre ») et non démontrés. Qu'il s'agisse du noyau de $F$ ou de l'image
de $V$, et de $M^{(p)}$ ou de $M$, dépend des conventions de la théorie
contravariante ; on reproduit celles de la page. Sur une base parfaite
(page 48), tout cristal de Dieudonné est admissible.} Une marge annonce
encore trois paragraphes qui ne sont pas écrits : passage à la limite sur un
anneau $p$-adiquement complet, théorème de Tate, dualité.

\pagerange{43}{47}
\subsection*{2) et 3) Schémas abéliens (pages 43 à 47)}

Pour les schémas abéliens il n'y a plus de nombre premier privilégié, mais
il faut le site cristallin \emph{nilpotent}. On pose
$\mathbb{D}^*(A) = R^1 f_{\mathrm{cris}*}(\mathcal{O}_{A_{\mathrm{cris}}})$, d'où
$\mathbb{D}^*(A)_S = H^1_{\mathrm{DR}}(A/S)$ et, pour tout relèvement $A'$ de $A$
à un épaississement $S'$ à puissances divisées localement nilpotentes,
$\mathbb{D}^*(A)_{S'} = H^1_{\mathrm{DR}}(A'/S')$. Un relèvement existe toujours
sur $S$ affine (il renvoie à la \emph{Geometric Invariant Theory} de
Mumford), et l'indépendance du choix est, dit-il, une conséquence du
« tapis » de la cohomologie cristalline. La filtration de Hodge
\[
  0 \to f_*\Omega^1_{A/S} = \omega_A \to H^1_{\mathrm{DR}}(A/S)
  \to R^1 f_*\mathcal{O}_A = \mathfrak{t}_{A^*} \to 0
\]
factorise $\mathbb{D}^*$ par des cristaux filtrés, et :

\textbf{Théorème de déformation pour les schémas abéliens.} Si $S'$ est un
épaississement de $S$ à puissances divisées localement nilpotentes, la donnée
d'un schéma abélien $A'$ sur $S'$ équivaut à celle de $A = A'|_S$ et d'un
relèvement localement facteur direct de $\omega_A$ dans $\mathbb{D}^*(A)_{S'}$.

Une seconde description passe par l'\emph{extension vectorielle universelle}
$0 \to \omega_{A^*} \to E(A) \to A \to 0$ : elle s'étend en un cristal en
groupes lisses $\underline{E}(A)$, avec $\underline{E}(A)_{S'} = E(A')$ pour tout
relèvement $A'$, et l'on a un isomorphisme canonique
$\mathbb{D}^*(A) \simeq \underline{\mathrm{Lie}}\,\underline{E}(A^*)$ compatible aux
filtrations. Relever $\omega_A$ en $\mathrm{Fil}^1$, c'est relever le
sous-groupe vectoriel correspondant $L$ de $\underline{E}(A^*)_{S'}$, et le
quotient $\underline{E}(A^*)_{S'}/L$ est le dual du schéma abélien $A'$
cherché. En caractéristique nulle, le relèvement de $E$ est unique à
isomorphisme unique près — « le résultat déjà signalé dans ma vieille lettre
à Tate » ; en général on procède par « la méthode de l'exponentielle »,
annoncée pour la suite du séminaire.\footnote{L'extension vectorielle
universelle et la description $H^1_{\mathrm{DR}}(A/S) \simeq \mathrm{Lie}\,E(A^*)$
sont exposées dans B.~Mazur et W.~Messing, \emph{Universal extensions and
one dimensional crystalline cohomology}, LNM 370 (1974). La « vieille lettre
à Tate » n'est pas identifiée ici.} Il serait plus naturel, note-t-il, de
travailler avec $\mathbb{D}_*(A) = \underline{\mathrm{Lie}}\,\underline{E}(A)
\simeq \mathbb{D}^*(A^*)$, les deux points de vue s'échangeant par
l'accouplement parfait $\mathbb{D}^*(A) \otimes \mathbb{D}^*(A^*) \to
\mathcal{O}$. Une remarque (page 46) annonce que la théorie des groupes de
Barsotti--Tate donnera de même deux constructions de $\mathbb{D}^*(G)$, l'une
cohomologique, l'autre par une extension vectorielle universelle
$\underline{E}(G)$, elle-même obtenue soit cohomologiquement, soit « par
l'exponentielle », la voie adaptée au théorème de déformation.

Enfin (page 47), si $p$ est localement nilpotent, on attend un isomorphisme
fonctoriel $\mathbb{D}^*(A[p^\infty]) \simeq R^1 f_{\mathrm{cris}*}
\mathcal{O}_{A_{\mathrm{cris}}}$ compatible aux filtrations — c'est, ajoute-t-il,
en postulant un tel isomorphisme qu'on arrive heuristiquement à la
définition de $\mathbb{D}^*(G)$. Il en résulte que les déformations de $A$ et
de $A[p^\infty]$ sont les mêmes, d'où le théorème de Serre--Tate (« méthode
suivie par Messing dans sa thèse »).

\pagerange{48}{53}
\subsection*{« Rabiots », problèmes, et un théorème de Tate (pages 48 à 53)}

Une liste de compléments, sur un autre papier, ajoute :

\emph{a)} Sur une base parfaite, la filtration canonique est caractérisée par
$\mathrm{Fil}^{1\,(p)} = \operatorname{Ker} F_M = \operatorname{Im} V_M$, et tout
cristal de Dieudonné est admissible.

\emph{b)} Les \emph{isogroupes} de Barsotti--Tate : on remplace
$\mathrm{Hom}(G,H)$ par $\mathrm{Hom}(G,H) \otimes \mathbb{Q}$ (définition
« provisoire », qui ne fournit pas un champ). \emph{Proposition} (« prouvée
plus tard ») : si $S \subset S'$ est une immersion nilpotente, $p$ localement
nilpotent, la restriction $\mathrm{IsoBT}(S') \to \mathrm{IsoBT}(S)$ est une
équivalence.\footnote{La page écrit « immersion nilpotente d'Idéal $J$ annulé
par une puissance de $J$ », phrase qui se répète ; on comprend $J$ nilpotent.
La pleine fidélité est le lemme de rigidité des quasi-isogénies (Drinfeld ;
voir N.~Katz, « Serre--Tate local moduli », LNM 868, 1981) ; nous ne nous
prononçons pas sur l'essentielle surjectivité telle qu'énoncée.} Il en tire
« l'énoncé VII 10.8.2 de mes notes de Pise », que nous n'identifions pas.

\emph{Problème 4} : décrire l'image essentielle de ce foncteur pour $A$
noethérien, séparé et complet pour la topologie $p$-adique ; y a-t-il un
critère valuatif ? Pour $A$ de valuation discrète d'inégales
caractéristiques, est-ce un ouvert de Zariski d'une grassmannienne sur le
corps des fractions $L$ ? Même pour $A = W(k)$, $k$ algébriquement clos, le
foncteur n'est pas essentiellement surjectif.

\emph{Problème 5} : pour $A$ de valuation discrète complet d'inégales
caractéristiques, à corps résiduel parfait $k$, expliciter le foncteur qui va
des couples $(M, \mathrm{Fil})$ — un module de Dieudonné sur $k$ et une
filtration à deux crans de $M \otimes_W L$ — vers les représentations
$p$-adiques de $\mathrm{Gal}(\overline{L}/L)$ ; peut-on le définir sans se borner
à l'image, en imposant seulement des conditions de rang ?

\emph{Problème 6} : la même question sur une base $A$ plus générale, vers les
faisceaux $p$-adiques sur $\mathrm{Spec}(A[1/p])$.\footnote{Une « Conjecture 7 »
commencée dans la même veine est biffée à la main. En tête de la page 52, au
crayon, deux inégalités $\sum_1^i \lambda'_j \leq \sum_1^i \lambda_j$,
$\sum_1^n \lambda'_j = \sum \lambda_j$ et « mieux de partir des Hodge » :
une comparaison de polygones, du type de celle que B.~Mazur démontrera entre
polygone de Newton et polygone de Hodge (1972-1973) ; rien sur la page ne dit
lequel est lequel.}

\textbf{Le théorème de Tate.} Soit $S$ normal, noethérien, intègre, de point
générique $\eta$ de caractéristique nulle. Alors $\mathrm{BT}(S) \to
\mathrm{BT}(\eta)$ est pleinement fidèle : un groupe de Barsotti--Tate sur
$S$ est connu, à isomorphisme unique près, par la représentation
$p$-adique de $\pi_1(\eta)$ qu'il définit.\footnote{J.~Tate, « $p$-divisible
groups », \emph{Proceedings of a Conference on Local Fields}, Driebergen
1966 (Springer, 1967) ; la page renvoie à « Driebergen ». Le dossier 100 est
bâti sur ce théorème (son tapuscrit des pages 47 à 50), et sa lecture
modernisée en discute la portée.} Pour $S$ le spectre d'un anneau de
valuation discrète complet $V$ d'inégales caractéristiques, à corps
résiduel parfait, deux questions se posent : caractériser les modules
galoisiens obtenus (Tate montre qu'ils sont de Hodge--Tate) ; trouver un
foncteur « explicite » de ces modules vers les F-cristaux sur $k$ munis d'une
filtration de $M \otimes_W L$. Le théorème suggère que le « \emph{foncteur
mystérieux} » des F-cristaux filtrés vers les représentations galoisiennes
est pleinement fidèle, et qu'il faut en expliciter un quasi-inverse sur une
sous-catégorie convenable, sans doute contenue dans celle des
représentations de Hodge--Tate.\footnote{C'est le nom sous lequel la question
est restée : J.-M.~Fontaine y a répondu par sa théorie des représentations
cristallines et des modules filtrés faiblement admissibles (avec
P.~Colmez, \emph{Invent. Math.} 140, 2000, pour « faiblement admissible
implique admissible »). Les feuillets ne contiennent aucune de ces notions.}

On conjecture enfin que l'hypothèse de caractéristique nulle sur $k(\eta)$ est
inutile ; le cas d'égales caractéristiques, dit-il, se ramènerait à la théorie
de Dieudonné (problème 1) et à la question suivante : pour $S$ normal
noethérien intègre de caractéristique $p$, la restriction
$\mathrm{Crisloclib}(S) \to \mathrm{Crisloclib}(\eta)$ (ou son analogue pour les
F-cristaux) est-elle pleinement fidèle ? On se ramène à $k[[t]]$, $k$
algébriquement clos, où la question porte sur des modules à connexion
formelle.\footnote{Le théorème de Tate en caractéristique $p$ a été démontré
par A.~J.~de~Jong (\emph{Invent. Math.} 134, 1998), par la voie cristalline ;
la lecture du dossier 100 le signale aussi.}

\pagerange{54}{59}
\section*{V. $\chi$ et $\chi_!$ (pages 54 à 67)}

Le titre est seul sur un feuillet (page 54) ; les pages suivantes sont à
l'encre bleue.

\subsection*{Les axiomes (pages 55 et 59)}

Soit $k$ un corps. On note $\mathcal{E}_{\chi_!}$ la classe des espaces
localement compacts $X$ dont la cohomologie à supports compacts
$H^*_!(X,k)$ est de dimension finie, et $\mathcal{E}_\chi$ celle des espaces dont
la cohomologie $H^*(X,k)$ l'est ; $\chi_!$ et $\chi$ sont les
caractéristiques d'Euler correspondantes.\footnote{Les marges disent
« cohomologie singulière » pour les deux. Pour $H^*_!$ des espaces localement
compacts, c'est la cohomologie des faisceaux (ou d'Alexander--Spanier) à
supports compacts qu'il faut entendre ; sur la classe des espaces localement
coniques où le dossier aboutit, les théories coïncident.} Il relève :
\begin{enumerate}
  \item[(a)] si $X$ est compact, $X \in \mathcal{E}_{\chi_!} \Leftrightarrow
  X \in \mathcal{E}_\chi$, et $\chi_!(X) = \chi(X)$ ;
  \item[(b)] si $Y \subset X$ est fermé, et deux des trois espaces $X$, $Y$,
  $X - Y$ sont dans $\mathcal{E}_{\chi_!}$, le troisième aussi, et
  $\chi_!(X) = \chi_!(Y) + \chi_!(X - Y)$ (suite exacte de cohomologie à
  supports compacts) ;
  \item[(c)] si $X = X' \cup X''$ avec $X'$, $X''$ fermés, et trois des quatre
  espaces $X, X', X'', X' \cap X''$ sont dans $\mathcal{E}_\chi$, le quatrième
  aussi, et $\chi(X) = \chi(X') + \chi(X'') - \chi(X' \cap X'')$
  (Mayer--Vietoris) ;\footnote{Pour un recouvrement fermé, la suite de
  Mayer--Vietoris vaut en cohomologie des faisceaux d'espaces paracompacts ;
  en cohomologie singulière il faut une hypothèse d'excision (paires de
  voisinages rétractiles), satisfaite dans le cadre localement conique. Une
  variante entre crochets (son étiquette « (d) » est biffée) note que
  $X, X' \cap X'' \in \mathcal{E}_\chi$ suffit à mettre $X'$ et $X''$ dans
  $\mathcal{E}_\chi$ : c'est ce que donne l'exactitude de la suite, le terme
  $H^*(X') \oplus H^*(X'')$ étant pris entre deux termes de dimension
  finie.}
  \item[(d)] si $X$ est dominé homotopiquement par $X'$ ($X \to X' \to X$
  homotope à l'identité), $X' \in \mathcal{E}_\chi \Rightarrow X \in
  \mathcal{E}_\chi$ ; en particulier $\chi$ est un invariant d'homotopie ;
  \item[(e)] un point est dans $\mathcal{E}_\chi$, $\chi(\mathrm{pt}) = 1$, et
  $\chi(\emptyset) = 0$ ;
  \item[(f)] $\chi_!(X \times Y) = \chi_!(X)\chi_!(Y)$ et
  $\chi(X \times Y) = \chi(X)\chi(Y)$ (formules de Künneth) ;
  \item[(g)] si $f : X \to Y$ est une fibration localement triviale de fibre
  $F$, avec $Y, F \in \mathcal{E}_{\chi_!}$, alors
  $\chi_!(X) = \chi_!(Y)\,\chi_!(F)$, et de même pour $\chi$.
\end{enumerate}
La démonstration de (g) passe par la suite spectrale de Leray
$H^p_!(Y, R^q f_!\,k) \Rightarrow H^{p+q}_!(X)$, et par un \emph{Lemme} : pour
un système local $\mathcal{R}$ de rang $b$ sur $Y$, $\chi_!(Y, \mathcal{R}) =
b\,\chi_!(Y)$. Il le démontre lorsque $Y$ admet une filtration finie par des
fermés dont les différences sont des sommes finies de cellules
$Z_{ij} \simeq \mathbb{R}^{d_{ij}}$, par additivité et parce qu'un système local
sur une cellule est constant.\footnote{La page écrit $X \in \mathcal{E}_\chi$
et $\chi(X)$ dans la démonstration de (g) là où l'énoncé porte sur
$\chi_!$ ; on rétablit $\chi_!$. Elle écrit aussi $\sum (-1)^p H^p(Y,
\mathcal{R})$ pour la somme alternée des \emph{dimensions} des $H^p_!$.}

\subsection*{Premières conclusions (page 57)}

Pour tout $n$, $\mathbb{R}^n$, $\overline{\mathbb{R}}^n = [-\infty,+\infty]^n$, la
boule $B^n$ et la sphère $S^n$ sont dans $\mathcal{E}_\chi \cap
\mathcal{E}_{\chi_!}$, et
\[
  \chi(\mathbb{R}^n) = \chi(\overline{\mathbb{R}}^n) = \chi(B^n) =
  \chi_!(\overline{\mathbb{R}}^n) = \chi_!(B^n) = 1,
\]
\[
  \chi(S^n) = \chi_!(S^n) = 1 + (-1)^n, \qquad
  \chi_!(\mathbb{R}^n) = (-1)^n .
\]
Les trois premiers sont contractiles, d'où $\chi$ par (d) et (e), et les
compacts relèvent de (a). Pour la sphère, $S^n = S^n_+ \cup S^n_-$ avec
$S^n_+ \cap S^n_- \simeq S^{n-1}$, d'où par (c) $\chi(S^n) = 2 - \chi(S^{n-1})$,
et par récurrence ($\chi(S^0) = 2$) $\chi(S^n) = 1 + (-1)^n$. Enfin
$\mathbb{R}^n \simeq S^n - \mathrm{pt}$, et (b) donne $\chi_!(\mathbb{R}^n) =
(1 + (-1)^n) - 1 = (-1)^n$.\footnote{La page écrit ici $\chi(\mathbb{R}^n)$ ;
c'est $\chi_!(\mathbb{R}^n)$ que (b) calcule, $\chi(\mathbb{R}^n)$ valant $1$.}

D'où la \emph{formule d'Euler} : si $X$ admet une filtration par des fermés
$\emptyset = X_{-1} \subset X_0 \subset \cdots \subset X_n = X$ telle que
$X_i - X_{i-1}$ soit réunion disjointe de $f_i$ cellules ouvertes
$\simeq \mathbb{R}^i$, alors $X \in \mathcal{E}_{\chi_!}$ et
$\chi_!(X) = \sum_i (-1)^i f_i$ — et si $X$ est compact c'est aussi
$\chi(X)$.\footnote{« Filtration de l'espace $X$ » : que les $X_i$ soient
fermés est sous-entendu ; (b) en a besoin. Un énoncé ① ajouté et encadré
donne l'additivité sur les sommes disjointes.}

\pagerange{61}{67}
\subsection*{Espaces localement coniques (pages 61 à 65)}

Le recto de la page 61 commence au milieu d'un argument, sur des conditions
$(\psi_1)$, $(\psi_2)$ qui ne figurent pas dans les pages conservées : il
s'agit de recouvrements finis arbitrairement fins par des fermés dont toutes
les intersections sont dans $\mathcal{E}_{\chi_!}$, et, pour $\chi$, d'un
lemme analogue au précédent en homologie.

Il définit alors $\mathcal{C}_{-1} = \{\emptyset\}$ et, pour $n \geq 0$,
$\mathcal{C}_n$ la classe des espaces compacts $X$ dont tout point $x$ a un
voisinage homéomorphe au cône sur un $S \in \mathcal{C}_{n-1}$, le sommet
allant en $x$ ; $\mathcal{C}_\infty = \bigcup_n \mathcal{C}_n$ est la classe
des compacts \emph{localement coniques}, et $\mathcal{E}_\infty$ celle des
espaces localement compacts séparés $X$ dont le compactifié d'Alexandroff
$\hat{X} = X \cup \{\omega\}$ est dans $\mathcal{C}_\infty$.\footnote{La page
écrit « $S \in \mathcal{E}_{n-1}$ » dans la définition de $\mathcal{C}_n$, et
sa marge énumère des « $\mathcal{E}_0$ … $\mathcal{E}_4$ » ; on lit partout
$\mathcal{C}$. La marge dit ces espaces localement triangulables en dimension
$\leq 4$ (« Moïse »), triangulables s'ils sont compacts en dimension $\leq 3$,
et ajoute que pour la dimension $4$ « il n'est pas clair qu'ils soient
triangulables ». Le théorème de Moise porte sur la dimension $3$ ; et en
dimension $4$ la réponse est négative : la variété $E_8$ de Freedman, compacte
et localement conique, n'est pas triangulable (Casson).}

\textbf{Proposition.} $\mathcal{C}_\infty \subset \mathcal{E}_\chi \cap
\mathcal{E}_{\chi_!}$, et $\mathcal{E}_\infty \subset \mathcal{E}_\chi \cap
\mathcal{E}_{\chi_!}$.

Un compact localement conique est localement contractile, donc dans
$\mathcal{E}_\chi$ (et dans $\mathcal{E}_{\chi_!}$ par (a)). Pour
$X \in \mathcal{E}_\infty$, (b) appliqué à $\hat{X}$ et au point $\omega$
donne $X \in \mathcal{E}_{\chi_!}$ et $\chi_!(X) = \chi(\hat X) - 1$. Pour
$\chi$, soit $C \simeq \mathrm{C}(S_\omega)$ un voisinage conique de $\omega$,
$T$ le cône homothétique de rapport $\tfrac12$, et $X' = \hat X - \mathring T$,
compact ; $X'$ est localement conique et $X'$ est homotopiquement équivalent
à $X$, d'où $X \in \mathcal{E}_\chi$. Par (c),
$\chi(X) = \chi(\hat{X}) - \chi(T) + \chi(\partial T)$ avec $\chi(T) = 1$ et
$\partial T \simeq S_\omega$, d'où
\[
  \chi(X) - \chi_!(X) = \chi(S_\omega),
\]
où $S_\omega$ est le « lien à l'infini » de $X$ : la différence des deux
caractéristiques « exprime la structure locale de $X$ à l'infini ».
(Pour $X = \mathbb{R}^n$ : $1 - (-1)^n = \chi(S^{n-1})$.)

\textbf{Théorème} (page 65). Sur $\mathcal{E}_\infty$, $\chi$ et $\chi_!$ sont
deux fonctions à valeurs entières satisfaisant (a) à (f), (g) au moins quand
la base est un polyèdre relatif, et
$\chi_!(X) = \chi(\hat X) - 1$, $\chi(X) - \chi_!(X) = \chi(S_\omega)$. Sur les
polyèdres relatifs (différences de deux polyèdres, l'un plongé linéairement
dans l'autre), elles sont caractérisées par (a) à (e) et données par la
formule d'Euler.\footnote{Deux marges disent ce qui n'est pas acquis :
« j'ignore si c'est vrai sans telle restriction » (pour (g)), et « j'ignore si
$\chi$ et $\chi_!$ sont caractérisées par a) à e) sur tout
$\mathcal{E}_\infty$ ». La page énonce (c) sans la condition $X = X' \cup X''$,
qu'on rétablit. Le théorème dit « deux fonctions uniques » ; l'unicité n'est
établie que sur les polyèdres relatifs. La caractéristique d'Euler à supports
compacts comme invariant additif, « combinatoire », des ensembles
constructibles a été systématisée depuis (calcul d'Euler sur les fonctions
constructibles, géométrie o-minimale) ; ces pages n'en ont pas le
vocabulaire.}

Une feuille isolée (page 67) porte une liste de thèmes : « orientations ;
plongements ; $H^i$, $\chi$ ; rev.\ et $\pi_1$ ; rel.\ entre $\pi_1$ et th.\ de
Galois ; rev.\ et fonctions de var.\ complexes ; Conj.\ de Weil ». Les deux
premières chemises du dossier en traitent les deux premiers points ; on n'en
infère pas davantage.

\pagerange{71}{73}
\section*{VI. Invariants différentiels des applications étales d'une courbe (pages 71 à 83)}

Changement de papier : des feuillets perforés à l'encre bleu-noir, dont les
marges gauches sont couvertes d'ajouts écrits en long, lus en partie
seulement. Le sujet est celui de la page 35, devenu théorème.

\subsection*{L'invariant projectif, et son unicité (pages 71 à 73)}

Soit $X$ une courbe analytique complexe régulière, $\mathcal{E}_X$ (resp.
$\mathcal{P}_X$) le faisceau des applications étales d'ouverts de $X$ dans
$\mathbb{C}$ (resp. dans $\mathbb{P}^1(\mathbb{C})$).

\textbf{Théorème I.} Il existe un morphisme de faisceaux
$G : \mathcal{E}_X \times \mathcal{E}_X \to \omega_X^{2}$ tel que :
\begin{enumerate}
  \item[a)] $G$ est un opérateur différentiel non linéaire : en termes d'une
  uniformisante locale $t$, $G(f,g) = R(f,\ldots,f^{(n)}, g, \ldots,
  g^{(n)})\,dt^2$ avec $R$ rationnelle, dont le dénominateur ne s'annule
  identiquement pour aucun couple de fonctions étales et dont la valeur n'a
  pas de pôle là où $f'$ et $g'$ n'ont pas de zéro ;
  \item[b)] $G(f,g) + G(g,h) = G(f,h)$ ;
  \item[c)] $G(f, ag) = G(f, g + b) = G(f,g)$ pour $a \in \mathbb{C}^*$,
  $b \in \mathbb{C}$, et $G(f, 1/g) = G(f,g)$ si $g$ n'a pas de zéro ;
  \item[d)] si $G(f,g) = 0$ et si $f$, $g$ coïncident à l'ordre $2$ en un point
  d'un ouvert connexe $U$, alors $f = g$ sur $U$ ;
  \item[e)] pour $f$ fixé, $g \mapsto G(f,g)$ est un épimorphisme de
  $\mathcal{E}_U$ sur $\omega_U^2$.
\end{enumerate}
Deux tels morphismes diffèrent par un facteur $\varphi \in
\Gamma(X, \mathcal{O}_X^*)$. \textbf{II.} Si l'on demande en outre la
compatibilité aux morphismes étales entre courbes variables, les
coefficients de $R$ sont des constantes, et l'indétermination se réduit à
une constante $c \in \mathbb{C}^*$. On peut prendre
\[
  G(f,g) = \Bigl(\Bigl[-2\,\frac{g'''}{g'} + 3\Bigl(\frac{g''}{g'}\Bigr)^2\Bigr]
  - \Bigl[-2\,\frac{f'''}{f'} + 3\Bigl(\frac{f''}{f'}\Bigr)^2\Bigr]\Bigr)dt^2 ,
\]
et l'on trouve un torseur canonique $\mathcal{T}_X$ sous $\omega_X^2$,
fonctoriel pour les morphismes étales, et une application
$c^X : \mathcal{P}_X \to \mathcal{T}_X$ telle que $G(f,g) = c^X(g) - c^X(f)$,
qui induit un isomorphisme $H\backslash\mathcal{P}_X \xrightarrow{\sim}
\mathcal{T}_X$, où $H = \mathrm{PGL}_2(\mathbb{C})$.\footnote{Il écrit
$\mathrm{GP}(1,\mathbb{C})$ pour $\mathrm{PGL}_2(\mathbb{C})$ et
$\mathbb{E}^1$ pour la droite affine. La condition c) n'énonce que
l'invariance par les trois types d'homographies ; ils engendrent le groupe
(page 101). Ces énoncés sont démontrés à la station IX ; ici ils sont
posés.}

\textbf{III. Le cas relatif.} Pour $X$ lisse de dimension relative $1$ sur un
schéma $S$ variable, avec compatibilité aux changements de base et aux
morphismes étales, les « dérivées » doivent être remplacées par les
opérateurs $D^{(i)}_t = \frac{1}{i!}\frac{d^i}{dt^i}$, et les coefficients
doivent être dans $\mathbb{Z}$. On trouve encore existence, et unicité
\emph{au signe près}. L'expression $G$ ne convient plus, à cause de ses
coefficients $2$ et $3$ ; il faut prendre $G_0 = \frac{1}{12}G$ :
\[
  G_0(f,g) = \Bigl[\Bigl(-\frac{D^{(3)}_t g}{D^{(1)}_t g}
  + \Bigl(\frac{D^{(2)}_t g}{D^{(1)}_t g}\Bigr)^2\Bigr)
  - \Bigl(-\frac{D^{(3)}_t f}{D^{(1)}_t f}
  + \Bigl(\frac{D^{(2)}_t f}{D^{(1)}_t f}\Bigr)^2\Bigr)\Bigr] dt^2 .
\]
Le signe est fixé par une valeur test : $G_0(t, t^3) = \frac{2}{3t^2}\,dt^2$,
qui a un sens en caractéristique $2$ — où d'ailleurs les deux choix
coïncident.\footnote{Il essaie d'abord $G_0(t,t^2)$ et écrit
« $\frac{1}{2t}\,dt^2$ (au lieu de $-\frac{1}{2t}\,dt^2$) » ; le calcul direct
donne $\frac{1}{4t^2}\,dt^2$, comme le note la transcription. Il écarte ce test
lui-même : si $2$ n'est pas inversible, $t^2$ n'est jamais une
uniformisante. Les deux valeurs sont vérifiées en Lean
(lean/Grothendieck/Folder93.lean, \texttt{N\_X\_sq} et \texttt{N\_X\_cube}),
sous forme débarrassée des dénominateurs : le numérateur
$(D^{(2)}s)^2 - D^{(3)}s\,D^{(1)}s$ vaut $1$ pour $s = t^2$ et $6t^2$ pour
$s = t^3$, d'où $G_0(t,t^2) = \frac{1}{4t^2}\,dt^2$ là où $2$ est inversible
et $G_0(t,t^3) = \frac{2}{3t^2}\,dt^2$ là où $3$ l'est.}

\emph{Remarques} (pages 72 et 73). 1) Le rôle de $\omega^2$ est universel : un
morphisme $G' : \mathcal{E}_X \times \mathcal{E}_X \to \mathcal{F}$ vers un faisceau
cohérent, ayant les mêmes propriétés, s'écrit de façon unique $G' = uG$ avec
$u : \omega^2 \to \mathcal{F}$, et $u$ est un isomorphisme. 2) Le torseur
$\mathcal{T}^G$ dépend du choix de $G$ : $\mathcal{T}^{\varphi G}$ est le
« multiple » de $\mathcal{T}^G$ par $\varphi$, et le choix le plus fin est
$\mathcal{T}^{G_0}$. 3) En caractéristique $p > 0$ donnée, on peut chercher à
caractériser $G_0$ par a), b), c) et la non-nullité ; l'indétermination
serait une constante de $\mathbb{F}_p^*$, ou de $\Gamma(S, \mathcal{O}_S^*)$ sur
une base fixée.\footnote{La fin de la page 73 est presque illisible ; on ne
garde que ce qui se lit.}

\pagerange{74}{83}
\subsection*{D'autres groupes (pages 74 à 83)}

Ce qui précède était centré sur $H = \mathrm{PGL}_2$ opérant sur
$\mathbb{P}^1$. Il passe en revue les analogues pour d'autres groupes opérant
sur un espace homogène de dimension $1$ — ce qui « semble condition de
succès », $\mathbb{C}^*$ opérant sur $\mathbb{C}$, qui n'est pas homogène,
étant « mauvais ».
\begin{enumerate}
  \item[a)] \emph{Translations sur la droite.} $d : \mathcal{O}_X \to \omega$ est
  surjectif de noyau les constantes ; on pose $A(f,g) = dg - df$. Il vaut
  mieux se restreindre aux fonctions étales, et remplacer $\omega$ par le
  sous-faisceau $\omega^*$ de ses bases (torseur sous $\mathcal{O}_X^*$),
  l'invariant devenant $dg/df \in \mathcal{O}_X^*$.
  \item[b)] \emph{Homothéties sur $\mathbb{G}_m$.} Pour un module inversible
  $L$, de faisceau des bases $L^*$, chaque base $f$ définit la connexion
  $D_f$ qui la rend horizontale ; d'où $L^*/\mathbb{C}^* \xrightarrow{\sim}
  \mathrm{Conn}(L)$, torseur sous $\omega$, et l'invariant
  $B(f,g) = d(g/f)/(g/f) = D_g - D_f$. Pour $L = \mathcal{O}_X$ c'est
  $g \mapsto dg/g$.\footnote{La page 75 annonce le cas $L = \omega$ et écrit
  l'exemple pour $\mathcal{O}_X^*$, c'est-à-dire $L = \mathcal{O}_X$ ; on garde
  l'exemple.}
  \item[c)] \emph{Groupe affine.} En composant a) et b) :
  $\mathrm{Aff}(1,\mathbb{C})\backslash\mathcal{E}_X \xrightarrow{\sim}
  \mathrm{Conn}(\omega)$ par $f \mapsto D_{df}$, la connexion de $\omega$ pour
  laquelle $df$ est horizontale, et l'invariant
  \[
    D(f,g) = \frac{d\,(dg/df)}{dg/df} = \Bigl(\frac{g''}{g'} -
    \frac{f''}{f'}\Bigr)dt .
  \]
  \item[d)] \emph{$\mathbb{G}_m$ et son automorphisme $f \mapsto 1/f$.} Sur
  $\mathcal{E}_X^*$ (applications étales vers $\mathbb{C}^*$), $f \mapsto
  (df/f)^2$, et l'invariant $E(f,g) = (dg/g)^2/(df/f)^2$ (page 83).
\end{enumerate}

\textbf{La table des « écarts » (page 81).} Une feuille dresse le tableau de
ces invariants — écart additif $dg/df$, sa puissance $n$-ième, écart affine
$D_{dg} - D_{df}$, écart multiplicatif $(dg/g)/(df/f)$, écart projectif
$G(f,g)$ — et les relie par une chaîne de quotients
\[
  \mathcal{E}_X \to \mathcal{E}_X/\mathbb{G}_a \simeq \omega^{*}
  \xrightarrow{\ \alpha_n\ } \mathcal{E}_X/(\mu_n\cdot\mathbb{G}_a) \simeq \omega^{n*}
  \xrightarrow{\ \beta_n\ } \mathcal{E}_X/\mathrm{Aff}_1 \simeq \mathrm{Conn}(\omega),
\]
\[
  \mathrm{Conn}(\omega) \xrightarrow{\ \gamma\ } \mathcal{E}_X/\mathrm{PGL}_2
  \simeq \mathrm{Conn}_X ,
\]
où $\alpha_n(\omega_1) = \omega_1^n$, $\beta_n(\psi) = D_{dt} - \frac1n \frac{d\psi}{\psi}$
pour une base $\psi$ de $\omega^n$ écrite relativement à $dt^n$, et où
$\gamma$ envoie une connexion $D$ de $\omega$ sur la connexion projective
$c_D$ qui la rend « horizontale ». La formule clé, encadrée, est
\[
  c_{D + \omega_1} = c_D + 2D\omega_1 + \omega_1^2 \qquad (\omega_1 \in \omega).
\]
Pour $\psi = f\,dt^n$ ($n \neq 0$), elle donne la connexion projective
attachée à une $n$-différentielle,
\[
  c_n(f\,dt^n) = c_t + \Bigl[-\frac{2}{n}\,\frac{f''}{f}
  + \frac{2n+1}{n^2}\Bigl(\frac{f'}{f}\Bigr)^2\Bigr]dt^2 ,
\]
qui pour $n = 1$ et $f = s'$ redonne $G(t,s)$.\footnote{C'est la formule de la
page 79, où le coefficient $\frac{2n+1}{n^2}$ est d'une lecture incertaine ; on
l'a vérifiée par le calcul, $c_{D_{dt} + \lambda\,dt} = c_t + (2\lambda' +
\lambda^2)\,dt^2$ avec $\lambda = -\frac1n f'/f$. La page 81 écrit $\beta_n$
avec le signe $+\frac1n$ ; avec ce signe le coefficient serait
$\frac{1-2n}{n^2}$ et le cas $n = 1$ ne redonnerait pas $G$ : on suit la
convention de la page 79 et de la flèche « $\psi \mapsto -\frac1n
\frac{d\psi}{\psi}$ » de la même page 81. Pour $n = 2$ le coefficient vaut
$\frac54$, que l'encadré de la page 81 porte, lu comme douteux. La même page
79 écrit $(g'/g)^2$ dans le second terme d'une différence où il faut
$(f'/f)^2$, et pose $c'_n = n^2 c_n/12$ dans un « $\mathcal{T}_X^{\otimes n^2}$ »
dont l'exposant est incertain ; on n'en tire rien.}

Les pages 79 et 80 sont des brouillons. Le second, traversé de deux traits,
porte un diagramme reliant revêtement universel et disques en un point
$x$, des produits tensoriels $\check{E} \otimes E$ pour $E = \Omega^{-1} \oplus
\Omega^0$ — la décomposition graduée en $\Omega^1, \Omega^0, \Omega^{-1}$ qui
sera celle de la station VII — et les nombres $g+1$, $3g-3$, $4g-2$ ; on ne
le restitue pas.

\textbf{Une page étrangère (page 82).} Une feuille isolée pose, pour un espace
de lacets $X = \Omega Y$, la question de savoir si le premier invariant de
Postnikov $k_1 \in H^3(\pi_1 X, \pi_2 X)$ est nul, en notant qu'il provient
de $k'_1 \in H^4(K(\pi_1 X, 2), \pi_2 X)$ (l'invariant de $Y$), et, après une
nouvelle délooping, d'un élément de $H^5(K(\pi_1 X, 3), \pi_2 X)$, déjà dans
le domaine stable $H^{5+i}(K(\pi_1X, 3+i), \pi_2X)$. La page ne répond pas ;
elle ne se rattache visiblement à rien de ce qui l'entoure.

\pagerange{85}{87}
\section*{VII. Fibrés en droites projectives et connexions transverses (pages 85 à 100)}

La station change d'échelle : au lieu de formules en coordonnée, des
structures sur un schéma $X$, puis sur une courbe $X/S$.

\subsection*{1.1 et 1.2. Données équivalentes (pages 85 et 86)}

Sur un schéma $X$, se correspondent :
\begin{enumerate}
  \item[a)] un fibré en droites projectives $f : P \to X$ (fibré de
  Brauer--Severi de rang $1$) ;
  \item[b)] un module localement libre $E$ de rang $2$, à tensorisation près
  par un module inversible ($P = \mathbb{P}(E)$, $E = f_*L$ pour $L$ de degré $1$
  sur les fibres) ;
  \item[c)] une algèbre d'Azumaya $M$ de rang $4$ ($M = \underline{\mathrm{End}}(E)$) ;
  \item[d)] une forme $\mathfrak{g}$ de $\mathfrak{pgl}_{2,X}$, et d$'$) une forme de
  $\mathfrak{sl}_{2,X}$ — ces deux-là si $2$ est inversible sur $X$ ;
  \item[e)] une forme de $\mathrm{PGL}_{2,X}$ ($G = \underline{\mathrm{Aut}}_X P$), et
  e$'$) une forme de $\mathrm{SL}_{2,X}$ ($G' = \underline{\mathrm{SL}}(E)$, son
  revêtement universel).
\end{enumerate}
On a $\mathfrak{g} = M/\mathcal{O}_X \simeq f_*(T_{P/X})$, avec le crochet des
champs de vecteurs verticaux, et $\mathfrak{g}' \subset M$ les éléments de trace
nulle ; $P$ est le schéma des sous-algèbres de Borel de $\mathfrak{g}$, ou des
sous-groupes de Borel de $G$. La marge avertit : « attention à la car.
(notamment car.~2 !) ».\footnote{C'est, à peu de chose près, la liste des
neuf catégories du « Dictionnaire lié à $\mathrm{GP}(1) \simeq \mathrm{SO}(3)$ »
des pages 57 à 59 du dossier 82, qui en fait la théorie en toute
caractéristique, et où les deux catégories d'algèbres de Lie sont, de même,
marquées comme n'étant équivalentes aux autres que si $2$ est inversible.
L'indice de « $\mathfrak{sl}(1)_X$ » est lu $\mathfrak{sl}(2)_X$.}

Si l'on se donne en outre une \emph{section} $\sigma$ de $P$ (1.2), on obtient
une extension $0 \to \Omega \to E \to \mathcal{O}_X \to 0$, $\Omega$ inversible ;
$P = \mathbb{P}(E)$ a pour section celle que définit le quotient
$\mathcal{O}_X$. Alors $\mathfrak{g} \simeq \mathrm{Sym}^2(E) \otimes \Omega^{-1}$
($\det E \simeq \Omega$) est filtré :
\[
  \mathrm{Fil}_1 = \Omega \subset \mathrm{Fil}_2 = \mathfrak{b} \subset
  \mathrm{Fil}_3 = \mathfrak{g}, \qquad
  \mathrm{gr}_1 = \Omega,\ \mathrm{gr}_2 = \mathcal{O}_X,\ \mathrm{gr}_3 = \Omega^{-1},
\]
où la sous-algèbre de Borel $\mathfrak{b}$ (stabilisateur de $\sigma$) est
isomorphe à $E$ comme extension de $\mathcal{O}_X$ par $\Omega$, et $\Omega$ est
son radical nilpotent.\footnote{La page écrit l'extension et la sous-algèbre
de Borel avec la même lettre $\mathcal{L}$, en établissant
$\mathrm{Fil}_2 \simeq \Omega^{-1} \otimes \operatorname{Ker}(\mathrm{Sym}^2 E \to
\mathcal{O}_X) \simeq E$. Elle note aussi que $\Omega$ et $\mathfrak{b}$ sont
orthogonaux l'un de l'autre pour la forme de Killing, et, à vérifier,
$\mathfrak{g}/\Omega \simeq \mathfrak{b}^\vee$. L'alinéa f) — une algèbre augmentée
nilpotente d'échelon $2$ — est lu de façon incertaine.} Deux sections
\emph{disjointes} $\sigma$, $\sigma'$ (1.3) équivalent à un module inversible
$\Omega$ ($P$ est la clôture projective du fibré vectoriel $W(\Omega)$, $\sigma$
la section à l'infini, $\sigma'$ la section nulle), ou à une graduation
$\mathfrak{g} \simeq \Omega \oplus \mathcal{O}_X \oplus \Omega^{-1}$, c'est-à-dire
au choix d'une sous-algèbre de Cartan et d'un des deux sous-groupes
unipotents qu'elle normalise.

\pagerange{88}{94}
\subsection*{1.4. Connexions, et la condition de transversalité (pages 88 à 94)}

Soit maintenant $X$ lisse sur $S$, et $\omega = \Omega^1_{X/S}$. On considère
des connexions (relatives à $S$, sans hypothèse de platitude) sur
$\mathfrak{g}$, compatibles au crochet — c'est-à-dire sur $P$. Localement,
après choix d'un scindage $\mathfrak{g} = \Omega \oplus \mathcal{O}_X \oplus
\Omega^{-1}$ et d'une base $\delta$ de $\Omega$, il y a une connexion d'origine
$c_0$ (qui annule $\delta$, $1$, $\delta^{-1}$), et toute autre s'écrit
$c = c_0 + \mathrm{ad}(\varpi)$, $\varpi = \varpi_2 + \varpi_1 + \varpi_0$ avec
$\varpi_2 \in \Omega \otimes \omega$, $\varpi_1 \in \omega$, $\varpi_0 \in
\Omega^{-1} \otimes \omega$ (toute dérivation de $\mathfrak{g}$ est
intérieure).

\textbf{Proposition.} Toute connexion $c$ sur $\mathfrak{g}$ vérifie
$c(\mathrm{Fil}_i\,\mathfrak{g}) \subset \mathrm{Fil}_{i+1}(\mathfrak{g}) \otimes
\omega$, et les applications induites
\[
  c_{1,2} : \mathrm{gr}_1 = \Omega \to \mathrm{gr}_2 \otimes \omega = \omega, \qquad
  c_{2,3} : \mathrm{gr}_2 = \mathcal{O}_X \to \mathrm{gr}_3 \otimes \omega =
  \Omega^{-1} \otimes \omega
\]
sont $\mathcal{O}_X$-linéaires, données par $\varpi_0$, avec
$c_{1,2} = -2\,c_{2,3}$.\footnote{La relation, encadrée, est indépendante des
scindages ; le facteur $-2$ dépend de la normalisation du crochet que la
page fixe page 87 ($[\omega_1, \omega_{-1}] = 2\,\omega_1\omega_{-1}$).}

La connexion respecte la filtration si et seulement si $\varpi_0 = 0$, ce
qui revient à dire qu'elle respecte la Borel, ou que $\sigma$ est
horizontale. À l'opposé, on dit que $c$ est \emph{transverse} à $\sigma$ — ou
que $\sigma$ est « nulle part stationnaire » — si $c_{2,3}$ est un
isomorphisme $\mathcal{O}_X \xrightarrow{\sim} \Omega^{-1} \otimes \omega$,
c'est-à-dire si $i = c_{2,3}$ identifie $\Omega$ à $\omega$ (pages 90 et
91).\footnote{L'interprétation de $c_{2,3}$ comme « dérivée » de la section
$\sigma$ relativement à la connexion passe par $\Omega \simeq
\sigma^*(\Omega^1_{P/X})$ (page 91) ; la phrase est en partie conjecturale
sur la transcription.}

\textbf{La forme canonique} (pages 90 à 93). Supposons $2$ inversible et $c$
transverse, $i : \Omega \xrightarrow{\sim} \omega$. La restriction
$c|_\Omega : \Omega \to \mathfrak{b} \otimes \omega$ est un opérateur
différentiel d'ordre $1$ ; il se factorise par l'opérateur universel
$d^1 : \Omega \to P^1_{X/S}(\Omega)$ et une application linéaire
$\gamma : P^1_{X/S}(\Omega) \to \mathfrak{b} \otimes \omega$, morphisme
d'extensions de $\Omega$ par $\Omega \otimes \omega$ qui induit $-2i$ sur les
quotients et l'identité sur les sous-modules ($\gamma(\lambda\delta \otimes df)
= \lambda\delta \otimes df$). Après multiplication par $-\frac12$ sur $\omega$,
on obtient un isomorphisme d'extensions
\[
  \mathfrak{b} \xrightarrow{\ \sim\ } P^1_{X/S}(\omega) \otimes \omega^{-1}
\]
induisant $-\frac12\,\mathrm{id}_\omega$ et $\mathrm{id}_{\mathcal{O}_X}$, et par lequel
$c|_\Omega$ devient l'opérateur différentiel universel
$d^1_{\omega} : \omega \to P^1_{X/S}(\omega)$.\footnote{La page dit la même
chose en termes de torseurs : celui que définit l'extension $\mathfrak{b}$ se
déduit de celui de $P^1_{X/S}(\omega) \otimes \omega^{-1}$ par
« multiplication par $-2$ ». Le diagramme (R) porte $-2\,\mathrm{id}$, le
diagramme (RR) $-\frac12\,\mathrm{id}_\omega$, tous deux cerclés.} Ainsi :

\textbf{Proposition.} Pour $X/S$ une courbe lisse et $2$ inversible, les
triples $(\mathfrak{g}, \mathfrak{b}, c|_\Omega)$ formés d'une connexion
transverse sont uniques à isomorphisme unique près : $E \simeq
P^1_{X/S}(\omega)\otimes\omega^{-1}$, et $c|_\Omega = d^1_\omega$. Les connexions
de $\mathfrak{g}$ qui prolongent $c|_\Omega$ forment un torseur sous
$\Omega^2 \otimes \omega \simeq \omega^{2}$ (pages 93 et 94).

Ce sont, dans le vocabulaire d'aujourd'hui, les \emph{$\mathfrak{sl}_2$-opers}
de la courbe, et ce torseur est celui des \emph{connexions projectives}.%
\footnote{Le nom d'« oper » est dû à A.~Beilinson et V.~Drinfeld (années
1990) ; la notion de connexion projective sur une surface de Riemann est
classique (É.~Cartan, R.~Gunning). Les pages ne nomment ni l'une ni
l'autre ; elles trouvent la condition de transversalité et la forme
canonique.}
La page 94 propose de faire le lien avec la description $(P, \sigma, c)$ et
avec l'isomorphisme $\mathcal{O}_{\hat P} \simeq P^\infty_{X/S}$ — ce que fera
la page 108.

\pagerange{95}{100}
\subsection*{1.5 et 1.6. Sections horizontales (pages 95 à 100)}

Les sections de $P$ disjointes de $\sigma$ sont les scindages de $E$,
c'est-à-dire (dans le cas canonique) les scindages de
$P^1_{X/S}(\omega)$, ou encore les \emph{connexions sur $\omega$} : les
opérateurs $D : \omega \to \omega^2$ tels que $D(f\eta) = fD\eta + \eta\,df$.
Elles forment un torseur sous $\omega$.

\textbf{Proposition} (page 95). Pour chaque scindage, il y a une et une seule
connexion transverse sur $\mathfrak{g}$ pour laquelle la section correspondante
de $P$ est horizontale.

En coordonnée $t$, prenant pour référence la connexion $c_0$ qui rend
horizontale la section $D_t$ (la connexion de $\omega$ pour laquelle $dt$ est
horizontale), une connexion projective s'écrit $c_0 + \mu\,dt^2$, une
connexion sur $\omega$ s'écrit $D_t + \lambda\,dt$, et la section
correspondante est horizontale si et seulement si
\[
  2\,\frac{d\lambda}{dt} + \lambda^2 = \mu ,
\]
équation de Riccati.\footnote{La page 96 écrit $\frac{d\lambda}{dt} + \lambda^2
= \varphi$, « sauf erreur de calcul », avec en marge « facteur 2 ? » ; la
page 102 retrouve la même équation et note « facteur 2 qui manque — erreur de
calcul ? ». Le facteur est rétabli page 105 et dans l'équation (1) de la page
113, et il découle de la formule $c_{D+\omega_1} = c_D + 2D\omega_1 +
\omega_1^2$ (pages 81 et 104) : on l'adopte partout. Par
$\lambda = 2u'/u$ l'équation devient l'équation linéaire $u'' =
\frac{\mu}{4}\,u$.} Sur une courbe lisse en caractéristique nulle, il y a au
plus une solution de valeur donnée en un point, et une solution formelle,
en général transcendante ; ses pôles sont les points où la section qu'elle
définit rencontre $\sigma$.

\textbf{Sections quelconques} (1.6). Une section de $P = \mathbb{P}(P^1_{X/S}(\omega))$
est un quotient inversible $P^1_{X/S}(\omega) \twoheadrightarrow L$,
c'est-à-dire un opérateur différentiel d'ordre $1$, $D : \omega \to L$, de
symbole $\sigma_0 \in \Gamma(L \otimes \omega^{-2})$, tel que
$D(f\eta) - fD(\eta) = \sigma_0(\eta\,df)$ ; elle rencontre la section
canonique aux zéros du symbole, les \emph{points de dégénérescence} de $D$.
Les sections non disjointes de $\sigma$ sont donc les connexions sur $\omega$
\emph{à singularités polaires} (page 100).

\textbf{Quand une telle section est-elle horizontale ?} (pages 97 à 99). Il
faut que l'ouvert de non-dégénérescence soit dense dans chaque fibre, et que
la connexion projective qu'elle détermine sur cet ouvert se prolonge en une
section régulière. Condition locale aux points de dégénérescence : si
$\lambda_t = b/a$ (dans $D(f\,dt) = a\,df/dt + bf$), $\lambda_t$ n'a que des
pôles simples, de partie principale
\[
  \lambda_t = \frac{2}{t} + (\text{pas de terme constant}) +
  O(t),
\]
car $2\lambda' + \lambda^2$ doit rester régulier là où $\lambda$ a un pôle.
\footnote{Avec $\lambda = a/t + c + O(t)$, la partie polaire de $2\lambda' +
\lambda^2$ est $(a^2 - 2a)/t^2 + 2ac/t$ : elle s'annule si et seulement si
$a = 2$ et $c = 0$. Les pages 99 et 106 disent exactement cela (le « 2 » de la
page 106 est écrit sur un « 1 ») ; la page 98 écrit « $\lambda^2 +
d\lambda/dt$ », qui donnerait le résidu $1$. Une partie des pages 97 à 99 est
pâle et serrée ; la condition « la section se prolonge en section
régulière » et la description par un « champ d'éléments de contact d'ordre
$2$ » le long de $\sigma$ sont lues par fragments.} Il y a donc, le long de
$\sigma$, un champ d'éléments de contact d'ordre $2$, et les sections qui
admettent une connexion projective les rendant horizontales sont celles qui
ont ce contact en leurs points de rencontre avec $\sigma$. La page 100, un
brouillon numéroté « 1 », pose encore la question des automorphismes de $P$
qui changent $c$ en $c + \varpi_2$.

\pagerange{101}{107}
\section*{VIII. Changement de coordonnée, et le dictionnaire (pages 101 à 114)}

Les pages 101 à 105 sont numérotées par lui 5, 2, 2 bis, 3, 4 ; dans son
ordre, le calcul commence page 102 et finit page 101.

\subsection*{La formule de transformation (pages 101 à 107)}

À toute connexion $D$ sur $\omega$ correspond la connexion projective
$c_D$, l'unique qui rende horizontale la section de $P$ définie par $D$, et
\[
  c_{D_0 + \omega_1} = c_{D_0} + 2D_0(\omega_1) + \omega_1^2 ,
\]
qu'il vérifie par transitivité (« OK », page 104). L'application
$\omega_1 \mapsto 2D_0(\omega_1) + \omega_1^2$, $\omega \to \omega^2$, n'est ni
linéaire ni semi-linéaire. Pour deux uniformisantes $t$, $s$, on a
$D_s = D_t + s'\,d(1/s') = D_t - (s''/s')\,dt$, d'où
\[
  c_{D_s} = c_{D_t} + \Bigl[-2\Bigl(\frac{s''}{s'}\Bigr)' +
  \Bigl(\frac{s''}{s'}\Bigr)^2\Bigr]dt^2
  = c_{D_t} + \frac{3s''^2 - 2s's'''}{s'^2}\,dt^2 = c_{D_t} + G(t,s).
\]
\textbf{Objection} (page 101) : $c_{D_s} = c_{D_t}$ si et seulement si
$3s''^2 = 2s's'''$, soit $s''^2 = C s'^3$, soit — en intégrant — $s$
\emph{fonction homographique} de $t$. Les trois types d'homographies
$t \mapsto t + a$, $t \mapsto bt$, $t \mapsto 1/t$ s'y vérifient
directement.

Conclusion (page 103) : « changer de connexion projective revient simplement
à changer le second membre de l'équation ». Une section rationnelle de $P$,
donnée par une forme $\omega_1$ à pôles, est horizontale pour une connexion
projective (alors unique) si et seulement si $2D_0(\omega_1) + \omega_1^2$
reste régulière aux pôles — d'où la condition « résidu $2$, pas de terme
constant » de la station VII (page 106).

\pagerange{108}{114}
\subsection*{Le dictionnaire (pages 107 à 114)}

\textbf{Structures équivalentes} (pages 108 et 109). Sur une courbe lisse
$X/S$, les données suivantes se correspondent :
\begin{enumerate}
  \item[(a)] un triple $(P, \sigma, c)$ : fibré en droites projectives, section,
  connexion relative à $S$ transverse à $\sigma$ ;
  \item[(a$'$)] $(P, \sigma)$ et un isomorphisme d'algèbres augmentées
  $\mathcal{O}_{(\hat P, \sigma)} \simeq P^\infty_{X/S}$ tel que la connexion
  transportée se globalise ;
  \item[(b)] une algèbre de Lie de type $\mathfrak{sl}_2$ quasi-scindée
  $\Omega \subset \mathfrak{b} \subset \mathfrak{g}$ munie d'une connexion transverse ;
  (b$'$) la forme canonique de la station VII ;
  \item[(c)] sur $\mathbb{P}_X = \mathbb{P}(P^1_{X/S}(\omega))$, avec l'isomorphisme
  canonique $\mathcal{O}_{\mathbb{P}_X}/\mathcal{J}^3 \simeq P^2_{X/S}$, un
  isomorphisme $\mathcal{O}_{\mathbb{P}_X}/\mathcal{J}^4 \simeq P^3_{X/S}$ qui le
  prolonge ;
  \item[(d)] une équation différentielle du troisième ordre sur
  $\mathbb{P}^1_X$, invariante par $\mathrm{PGL}_{2,X}$ ;
  \item[(e)] ($S = \mathrm{Spec}\,\mathbb{C}$) un faisceau d'applications
  analytiques étales de $X^{\mathrm{an}}$ dans $\mathbb{P}^1(\mathbb{C})$ qui soit un
  torseur sous $\mathrm{PGL}_2(\mathbb{C})$ — une \emph{structure projective} au
  sens des atlas.
\end{enumerate}
Il passe de (e) à (a) par le fibré associé au torseur et une section
canonique (les couples $(f, f(x))$, $f$ germe du torseur, sont tous
conjugués), et GAGA ; de (a) à (e) en trivialisant $P$ par la connexion au
voisinage d'un point et en lisant la section comme une application vers
$\mathbb{P}^1$.\footnote{(c) est la description par les jets d'ordre $3$ qui
reviendra page 133 : se donner une connexion projective, c'est relever à
l'ordre $3$ l'isomorphisme canonique d'ordre $2$. Les pages 109 à 112 sont
lues avec de nombreuses incertitudes ; le détail du passage (e) $\to$ (a), et
la description de $P$ comme $\mathcal{T}/B$ ($\mathcal{T}$ le torseur, $B$ un
Borel), sont restitués ici dans leurs grandes lignes seulement.}

\textbf{Les équations explicites} (page 113). Dans une coordonnée $t$, avec
$c_0 = c_{D_t} + \mu\,dt^2$, et $P = \mathcal{T}/B$, $B_u \subset B$ le
radical unipotent :
\begin{enumerate}
  \item[a)] les sections horizontales de $P$ disjointes de $\sigma$ sont les
  connexions $D = D_t + \lambda\,dt$ avec $c_D = c_0$ : $2\lambda' + \lambda^2 = \mu$ ;
  \item[b)] les sections horizontales de $\mathcal{T}/B_u$ au-dessus d'elles
  sont les couples $(D, \varphi)$, $\varphi = \alpha\,dt$ base de $\omega$ avec
  $D\varphi = 0$ : $\alpha'/\alpha = -\lambda$ ;
  \item[c)] celles de $\mathcal{T}$ sont les uniformisantes $s$ telles que
  $D_s = D$ : $s''/s' = -\lambda$, et $\varphi = ds$, $\alpha = s'$ ;
\end{enumerate}
et l'élimination de $\lambda = -s''/s'$ donne
\[
  G_t(s) = -2\,\frac{s'''}{s'} + 3\Bigl(\frac{s''}{s'}\Bigr)^2 = \mu(t) .
\]
L'équation de Riccati, l'équation linéaire du premier ordre et l'équation
schwarzienne sont ainsi les trois étages d'un même torseur
$\mathcal{T} \to \mathcal{T}/B_u \to \mathcal{T}/B$.

La propriété essentielle de $G_t$ est son invariance par
$s \mapsto (as + b)/(cs + d)$. Par changement d'uniformisante $\tau$,
\[
  G_\tau(s)\,d\tau^2 = G_t(s)\,dt^2 - G_t(\tau)\,dt^2 .
\]
\footnote{La page écrit $G_\tau(s) = G_t(s)\,\tau'^2$, après un signe surchargé,
et la transcription ne la corrige pas. La formule juste est celle du texte,
qui découle de la relation de cocycle de la station IX ; elle se réduit à
une simple multiplication par un facteur seulement si $\tau$ est une
homographie de $t$.}

\textbf{Théorème} (page 113). Soit $X$ une courbe algébrique complète sur
$\mathbb{C}$ et $t : U \to \mathbb{C}$ une uniformisante sur un ouvert. Il existe
une unique section $\mu$ de $\omega_U^2$ telle que les solutions de
$G_t(s) = \mu$ soient exactement les restrictions des fonctions
uniformisantes de $X^{\mathrm{an}}$ — les inverses locaux de la projection du
revêtement universel, plongé dans $\mathbb{P}^1(\mathbb{C})$.

\emph{Question.} Cette section $\mu$, analytique, est-elle algébrique ? « Premier
cas douteux : $X = \mathbb{P}^1_{\mathbb{C}}$ moins trois points, ou une courbe
elliptique moins un point. »\footnote{Le premier n'est pas douteux, et le
dossier le règle lui-même à la page 137 : sur la droite projective moins
trois points il n'y a qu'une connexion projective régulière à l'infini, qui
est donc la connexion uniformisante. Le second l'est : c'est le premier cas
du « problème des paramètres accessoires » de Klein et Poincaré, ouvert au
sens où l'on ne sait pas exprimer ces paramètres algébriquement.} \emph{Autre
question} : quand $X$ varie dans une famille algébrique, la section $\mu$
dépend-elle algébriquement du paramètre, et sinon quelle est sa nature
« transcendante » ?\footnote{Sur l'espace de Teichmüller, la section
fuchsienne du fibré affine des structures projectives n'est pas même
holomorphe ; sa dérivée $\bar\partial$ est un multiple de la forme de
Weil--Petersson (P.~Zograf et L.~Takhtajan, 1987). La question appelle donc
une réponse négative.}

\pagerange{115}{120}
\section*{IX. Mise au net : le cocycle schwarzien (pages 115 à 123)}

D'une écriture soignée, la théorie est reprise depuis le début.

\textbf{Le cas analytique.} Soit $\mathcal{X}$ une courbe analytique et $t$, $s$
deux uniformisantes. On a :
\begin{enumerate}
  \item[(1)] $G(t, s + b) = G(t, as) = G(t, 1/s) = G(t,s)$ ($a \in
  \mathbb{C}^*$, $b \in \mathbb{C}$, $s$ sans zéro dans le troisième cas) ;
  \item[(2)] $G(t,s) + G(s,r) = G(t,r)$ ;
  \item[(3)] la formation de $G$ commute à la restriction aux ouverts.
\end{enumerate}
D'où $G(t,t) = 0$ et $G(s,t) = -G(t,s)$, l'invariance en la première variable
(1 bis), et l'extension de $G$ aux applications étales vers
$\mathbb{P}^1(\mathbb{C})$ (fonctions méromorphes à pôles simples, uniformisantes
ailleurs), avec $G(g\cdot t, g'\cdot s) = G(t,s)$ pour $g, g' \in H =
\mathrm{Aut}\,\mathbb{P}^1(\mathbb{C}) = \mathrm{PGL}_2(\mathbb{C})$. On obtient
$G : \mathcal{P}_X \times \mathcal{P}_X \to \omega_X^2$ invariant par $H \times H$
(I) et vérifiant la relation de cocycle (II). Le théorème d'existence et
d'unicité des équations différentielles donne la réciproque (I$'$) : sur un
ouvert connexe, $G(t,s) = 0$ entraîne $t = g \cdot s$ pour un $g \in H$ ; et,
$t$ fixé, $s \mapsto G(t,s)$ induit un isomorphisme
$H\backslash\mathcal{P}_X \xrightarrow{\sim} \omega_X^2$.

\textbf{Proposition} (page 117, « presque formelle à partir de (II) »). Il
existe, de façon unique, un torseur $\mathrm{Conn}_{\mathcal{X}}$ sous
$\omega^2_{\mathcal{X}}$ et un morphisme $c^X : \mathcal{P}_X \to
\mathrm{Conn}_{\mathcal{X}}$, $t \mapsto c_t$, tels que $G(t,s) = c_s - c_t$.
$\mathrm{Conn}_{\mathcal{X}}$ est le faisceau des \emph{connexions projectives}
de $\mathcal{X}$ ; $c^X$ est invariant par $H$ et induit
$H\backslash\mathcal{P}_X \xrightarrow{\sim} \mathrm{Conn}_{\mathcal{X}}$ ; tout est
compatible aux morphismes étales $\mathcal{X}' \to \mathcal{X}$.

Sur $\mathbb{P}^1(\mathbb{C})$ il y a une connexion projective canonique
$\mathcal{C} = c_{\mathrm{id}}$, et c'est la seule, car
$H^0(\mathbb{P}^1, \omega^2) = 0$ ($\omega^2$ est de degré $-4$). Pour toute
uniformisante, $c_t = t^*(\mathcal{C})$.

\textbf{Extension à la géométrie algébrique} (pages 118 à 120). Pour $X$ lisse
de dimension relative $1$ sur un schéma $S$ (ou formel), les mêmes formules
définissent $G(t,s)$ pour $t$, $s : X \to \mathbb{A}^1_S$ étales, puis pour des
applications étales vers $\mathbb{P}^1_S$, invariant sous
$H_S = \underline{\mathrm{Aut}}_S(\mathbb{P}^1_S)$ et vérifiant (II) ; d'où le faisceau
étale $\mathrm{Conn}_{X/S}$, torseur sous $\omega_{X/S}^2$, et
$c^X : \mathcal{P}_{X/S} \to \mathrm{Conn}_{X/S}$ avec $c^X(t) = t^*(\mathcal{C})$.
Plus généralement, tout fibré en droites projectives $\mathcal{X}_0$ sur $S$
porte une unique connexion projective, et $G(t,s) = s^*c_{\mathcal{X}_0} -
t^*c_{\mathcal{X}_0}$ pour $t, s : X \to \mathcal{X}_0$ étales.\footnote{« que le
lecteur (sic) écrira… ». Une marge, presque illisible, mentionne une condition
« $\geq 3$ » sur la caractéristique et $P^3_{X/S}(\omega)$.}

En caractéristique nulle, pour $X/S$ à fibres géométriquement connexes,
$G(t,s) = 0$ entraîne encore $s = g \circ t$ ; $f^{-1}(H_S)\backslash\mathcal{P}_X
\to \mathrm{Conn}_{X/S}$ est injectif, mais n'est plus surjectif. Une connexion
projective dans l'image est dite « intégrable par fonctions algébriques ».

En caractéristique $p > 0$ l'injectivité tombe aussi : les homographies à
coefficients dans $\mathcal{O}_{X'}$ (fonctions $p$-ièmes puissances relatives)
laissent $G$ invariant, et le mieux qu'on puisse espérer est
$F_X^{-1}(H_{X'})\backslash\mathcal{P}_X \hookrightarrow \mathrm{Conn}_{X/S}$.

\textbf{Le facteur $12$.} Avec les dérivées divisées $D^{(n)}_t = \frac{1}{n!}
\frac{d^n}{dt^n}$, qui ont un sens dès que $X/S$ est différentiellement lisse,
\[
  G(t,s) = 12\,G_0(t,s), \qquad
  G_0(t,s) = -\frac{D^{(3)}_t s}{D^{(1)}_t s} + \Bigl(\frac{D^{(2)}_t
  s}{D^{(1)}_t s}\Bigr)^2 .
\]
La page affirme la formule sans écrire le calcul.\footnote{Vérifié en Lean
(lean/Grothendieck/Folder93.lean, \texttt{G\_eq}) : $G = 12\,G_0$ est une
identité sur $\mathbf{Z}$, valable sur tout anneau commutatif, sans division
par $2$ ni par $3$ ; elle découle de $k!\,D^{(k)} = d^k/dt^k$, qui est
\texttt{Polynomial.factorial\_smul\_hasseDeriv} dans mathlib (les $D^{(k)}$
sont les dérivées de Hasse, \texttt{Polynomial.hasseDeriv}). La même
vérification, dans le modèle local d'un anneau de polynômes, établit la
relation de cocycle (II) pour $G_0$ (\texttt{N\_comp}, \texttt{G0\_comp}) :
écrite sans dénominateur, c'est une identité entre polynômes des
coefficients, sans hypothèse d'inversibilité, donc valable sur toute base, en
caractéristique $2$ et $3$ comprises ; et l'invariance (1) par les
homographies dont les coefficients sont annulés par $D^{(1)}$, $D^{(2)}$,
$D^{(3)}$ (\texttt{G0\_homography}).} Donc, si $12$ est nul sur $S$, toutes les $c_t$ sont égales et il y a une
connexion projective canonique ; et en toute caractéristique,
$\mathrm{Conn}_{X/S}$ se déduit, par le changement de groupe $12\,\mathrm{id} :
\omega^2 \to \omega^2$, d'un torseur plus fin $\mathrm{Conn}_{0,X/S}$ muni de
$c_0^X$ avec $c_0^X(s) - c_0^X(t) = G_0(t,s)$, et l'on écrira $c_0 =
\frac{1}{12}c$ (page 121).\footnote{De l'édition : dans les deux travaux récents
lus pour situer ce $G_0$, il ne se trouve pas. Chez Y.~Hoshi (« A note on
Schwarzian derivatives and Sugiyama--Yasuda locally exact differentials »,
Osaka J. Math. 61, 2024, § 3 et § 6), l'invariant projectif de niveau $1$ est
la dérivée schwarzienne classique, à coefficients $2$ et $3$, qui s'annule
identiquement pour $p \leq 3$ (remarque 6.4.2) ; l'invariant de niveau $2$ en
caractéristique $2$ y est celui de Sugiyama et Yasuda, à valeurs dans les
différentielles localement exactes et non dans $\omega^{\otimes 2}$ ; rien n'y
porte sur $p = 3$ au niveau $2$. Chez Y.~Wakabayashi (« Frobenius-Ehresmann
structures and Cartan geometries in positive characteristic »,
arXiv:2109.02826, 2021), $p$ est impair et aucune schwarzienne n'est écrite.
Les sources qui trancheraient (Ihara 1974, Hoshi 2020, Sugiyama--Yasuda 2020)
n'ont pas été lues ; on n'affirme donc rien sur la priorité.}

\pagerange{121}{123}
\subsection*{L'injectivité en caractéristique $p$ (pages 121 à 123)}

On peut espérer que, $\mathcal{O}'_X$ désignant le sous-faisceau des fonctions
annulées par $D^{(1)}, D^{(2)}, D^{(3)}$ — c'est $\mathcal{O}_X^{p}$ si $p \neq
2,3$, et $\mathcal{O}_X^{p^2}$ si $p = 2$ ou $3$ —, le quotient de $\mathcal{P}_{X/S}$
par les homographies à coefficients dans $\mathcal{O}'_X$ s'injecte dans
$\mathrm{Conn}^0_{X/S}$.\footnote{La description de $\mathcal{O}'_X$, en
marge de la page 121 et marquée « OK », est vérifiée en Lean dans le modèle
local d'un anneau de polynômes en caractéristique $p$
(lean/Grothendieck/Folder93.lean, \texttt{killed\_iff\_of\_gt},
\texttt{killed\_iff\_two}, \texttt{killed\_iff\_three}) : les polynômes
annulés par $D^{(1)}$, $D^{(2)}$, $D^{(3)}$ sont ceux en $t^p$ si $p > 3$, en
$t^4$ si $p = 2$, en $t^9$ si $p = 3$ ; l'argument est le théorème de Lucas,
\texttt{Choose.choose\_modEq\_choose\_mod\_mul\_choose\_div\_nat} dans mathlib.
Le niveau $p^2$ est nécessaire (\texttt{char\_two}, \texttt{char\_three}) : en
caractéristique $2$, $D^{(1)}(t^2) = 0$ mais $G_0(t, t + t^2) = dt^2$ ; en
caractéristique $3$, $D^{(1)}(t^3) = D^{(2)}(t^3) = 0$ mais
$G_0(t, t + t^3) = -dt^2$. La forme générale de ce critère (polynômes annulés
par les $D^{(k)}$, $0 < k < p^e$) est proposée à mathlib, PR ouverte
n° 44736, avec le lemme de divisibilité tiré de Lucas qu'elle emploie, PR
ouverte n° 44735.} Un cas local le confirme :

\textbf{Proposition.} Soient $A$ un anneau de caractéristique $p > 3$ avec
$\mathrm{Pic}(A) = 0$, $B = A[T]/(T^p - \alpha)$, $t$ la classe de $T$, et
$s \in B$ avec $D^{(1)}_t s \in B^*$. Si $c_0(s) = 0$, c'est-à-dire
$(D^{(2)} s)^2 = D^{(1)}s \cdot D^{(3)}s$, il existe $a, b, c, d \in A$, avec
$ad - bc \in A^*$, tels que $s = (at + b)/(ct + d)$.\footnote{La page ajoute
« avec $ad - bc = 1$ » ; c'est faux en général, et nous l'avons retiré. Si
une homographie de déterminant $1$ vaut $ut$ modulo $t^3$, alors $u$ est un
carré ; or dans $\mathbf{F}_5[T]/(T^5)$, $s = 2t$ vérifie $c_0(s) = 0$ et
$D^{(1)}s = 2 \in \mathbf{F}_5^*$, et $2$ n'est pas un carré modulo $5$ :
vérifié en Lean (lean/Grothendieck/Folder93.lean,
\texttt{det\_one\_counterexample}). Le déterminant est une unité, carré
seulement si le coefficient de $t$ dans $s$ l'est ; l'énoncé juste porte sur
$\mathrm{PGL}_2(A)$.}

\emph{Démonstration.} On se ramène à $\alpha = 0$ ; $D^{(1)}s \in B^*$ veut dire
$a_1 \in A^*$ ($s = \sum a_i t^i$), et, quitte à composer par des
homographies, $s = t + a t^d + \cdots$ avec $3 \leq d \leq p - 1$. Modulo
$t^{d-2}$, on trouve $D^{(3)}s\,D^{(1)}s \equiv a\binom{d}{3}t^{d-3}$ et
$(D^{(2)}s)^2 \equiv 0$, donc $a\binom d3 = 0$, et $\binom d3$ est inversible
pour $3 \leq d \leq p-1$, $p > 3$ : $a = 0$, et par récurrence $s = t$.
\emph{Cas $p = 2, 3$.} Il faut prendre $B = A[T]/(T^{p^2} - \alpha)$, et
$3 \leq d \leq p^2 - 1$. Par le théorème de Lucas, $\binom d3$ est non nul
modulo $p$ si et seulement si l'addition $3 + (d-3)$ en base $p$ se fait sans
retenue : pour $p = 2$ seul $d = 3$ intervient ($\binom33 = 1$) ; pour $p = 3$,
$3 = (0,1)$ en base $3$ et $d - 3 \leq 5$ a un chiffre des « $3$ » au plus égal
à $1$, donc pas de retenue — « on gagne encore ! ».\footnote{Un premier essai
pour $p = 3$, faux, est barré de trois diagonales ; la colonne de droite,
ajoutée après coup, donne l'argument retenu. La page dit « $A$ local, plus
généralement $\mathrm{Pic}(A) = 0$ », hypothèse qui sert à normaliser les
coefficients de l'homographie. La vérification en Lean
(lean/Grothendieck/Folder93.lean, \texttt{rigidity\_of\_gt},
\texttt{rigidity\_four}, \texttt{rigidity\_nine}) établit la Proposition
pour $\alpha = 0$, dans $A[T]/(T^p)$ pour $p > 3$ et dans $A[T]/(T^{p^2})$
pour $p = 2, 3$, sur tout anneau de caractéristique $p$ et sans hypothèse sur
$\mathrm{Pic}(A)$ : on peut prendre $d = 1$, et le déterminant est alors le
coefficient de $t$ dans $s$ ; pour $T^4$ ($p = 2$) l'énoncé vaut même sur
tout anneau, seul $\binom33 = 1$ intervenant. L'hypothèse sur
$\mathrm{Pic}(A)$ ne peut donc servir qu'à la réduction de $\alpha$ à $0$,
que la page fait sans argument et qui n'est pas vérifiée. Pour $p > 3$,
\texttt{rigidity\_iff\_of\_gt} donne aussi la réciproque : une homographie à
coefficients dans $A$, de dénominateur $ct + d$ avec $d \in A^*$, vérifie
$c_0 = 0$. La démonstration vérifiée suit les coefficients de
$s\,(ct + 1)$ au lieu de composer $s$ par des homographies, et le
coefficient $a\binom d3$ de la page y reparaît tel quel.}

\pagerange{125}{129}
\section*{X. Uniformisation et connexion projective canonique (pages 125 à 139)}

\subsection*{Domaines homogènes et revêtement universel (pages 125 à 129)}

Soit $P$ une droite projective complexe. Un ouvert non vide $D \subset
P^{\mathrm{an}}$ est \emph{homogène} si un sous-groupe de $\mathrm{Aut}(P)$ y agit
transitivement.

\textbf{Théorème 1} (classique). a) À conjugaison près, il y a exactement trois
domaines homogènes : $P$, $P$ moins un point, un disque. b) Le stabilisateur
$G$ de $D$ dans $\mathrm{Aut}(P)$ s'identifie à $\mathrm{Aut}(D)$, et $D$ est connexe
et simplement connexe. c) $(P, D) \mapsto D$ est une équivalence de la
catégorie des couples (droite projective, domaine homogène) sur celle des
courbes analytiques connexes et simplement connexes.

Les stabilisateurs sont $\mathrm{PGL}_2(\mathbb{C})$, $\mathrm{Aff}(1,\mathbb{C})$ et,
pour un disque, $\mathrm{PSL}_2(\mathbb{R})$.\footnote{La page écrit
$\mathrm{GP}(1,\mathbb{R}) \simeq \mathrm{GL}(2,\mathbb{R})/\mathbb{R}^* \simeq
\mathrm{SL}(2,\mathbb{R})/\mu_2(\mathbb{R})$. Le dernier isomorphisme est faux :
$\mathrm{PGL}_2(\mathbb{R})$ a deux composantes connexes, et
$\mathrm{PSL}_2(\mathbb{R})$ en est la composante neutre. C'est
$\mathrm{PSL}_2(\mathbb{R})$ qui stabilise le disque ; $\mathrm{PGL}_2(\mathbb{R})$
stabilise le cercle et échange le disque et son « polaire » $P - \overline{D}$.}
Il y a correspondance entre les formes réelles de $P$ ayant des points
réels (un « cercle » $P_{\mathbb{R}}(\mathbb{R})$), les formes réelles
correspondantes de $\mathrm{Lie}\,\mathrm{Aut}\,P = H^0(P, T_P)$, et les domaines
homogènes du troisième type à polarité près (les deux composantes du
complémentaire du cercle).\footnote{La page parle de « descente du corps de
base de $\mathbb{C}$ à $\mathbb{R}$ » dans le groupe $\mathrm{Aut}_{\mathbb{C}}P$ ; la
forme réelle compacte $\mathrm{SO}(3)$, qui correspond à une conique sans
point réel, ne fournit pas de disque. La correspondance vaut pour les formes
déployées, c'est-à-dire pour les structures réelles de $P$ elle-même.}

\textbf{Scholie.} Il y a un foncteur $\mathbb{P}$ des courbes analytiques
simplement connexes vers les droites projectives, avec un plongement
fonctoriel $X \hookrightarrow \mathbb{P}(X)(\mathbb{C})$ en domaine homogène,
unique à isomorphisme unique près. C'est le théorème d'uniformisation de
Koebe et Poincaré, mis sous forme fonctorielle.

Appliqué aux revêtements universels $\widetilde{X}(x)$ d'une courbe $X$, qui
forment un système local sur le groupoïde fondamental, il donne :

\textbf{Théorème 2.} Pour $X$ une courbe analytique connexe non singulière,
il existe, à isomorphisme unique près : un fibré en droites projectives
$\mathbb{P}_X$ sur $X$, une connexion intégrable sur $\mathbb{P}_X$, une section
$\sigma$, un ouvert $\widetilde{X} \supset \sigma(X)$ de $\mathbb{P}_X$ et un
morphisme $p : \widetilde{X} \to X \times X$ au-dessus de $X$, tels que
$p\sigma = \Delta_X$, que chaque $p_x : \widetilde{X}_x \to X$ soit un revêtement
universel, que $p$ soit compatible aux connexions, et que chaque
$\widetilde{X}_x$ soit un domaine homogène de $(\mathbb{P}_X)_x$.\footnote{Quatre
notes marginales, reliées par des accolades aux données et aux conditions,
semblent dire lesquelles déterminent les autres ; elles ne se lisent pas.}

Selon le type du revêtement universel : si c'est $\mathbb{P}^1$, $X =
\mathbb{P}^1$ et tout est trivial ; si c'est $\mathbb{C}$, $X$ est $\mathbb{C}$,
$\mathbb{C}^*$ ou une courbe elliptique (quotients par des réseaux de rang
$0$, $1$, $2$), $\widetilde{X} = \mathbb{P}_X - \sigma'$ pour une section
horizontale $\sigma'$, et $\widetilde{X}$ s'identifie au fibré tangent, la
connexion étant celle d'une structure \emph{affine}. La question de son
algébricité ne se pose que pour $\mathbb{C}$ et $\mathbb{C}^*$, le cas compact
étant « réglé par GAGA » ; pour $\mathbb{C}$ tout est algébrique.

\pagerange{130}{137}
\subsection*{Le torseur des connexions projectives, et les singularités régulières (pages 130 à 137)}

On retient de ces structures celles qui ont un sens algébrique :
$\mathbb{P}_X$, la connexion $c_X$, la section $\sigma$, et l'isomorphisme
formel $\widehat{\mathbb{P}}_X \simeq \widehat{X \times X}$ le long de $\sigma$
et de la diagonale ; la condition est que $\sigma$ soit nulle part
stationnaire. « Les travaux de Deligne » donnent, si $X$ n'est pas propre,
l'existence et l'unicité de structures \emph{algébriques} induisant les
structures transcendantes — l'unicité, pas l'existence, de $\sigma$ — sous la
condition que la connexion soit \emph{régulière à l'infini}.\footnote{P.~Deligne,
\emph{Équations différentielles à points singuliers réguliers}, LNM 163
(1970) : la correspondance de Riemann--Hilbert algébrique. La page ne donne
pas la référence.}

\textbf{Le torseur algébrique} (pages 132 à 134). Pour une courbe lisse $X/S$
en caractéristique nulle, la catégorie des triples $(P, \sigma, c)$
transverses est rigide : $P = \mathbb{P}(E)$ avec $E = \mathcal{J}/\mathcal{J}^3 \otimes
\omega^{-1}$, extension de $\mathcal{O}_X$ par $\omega$ ($\mathcal{J}$ l'idéal
d'augmentation de $P^\infty_{X/S}$), $\sigma$ définie par cette extension, et
$c$ déduite d'un isomorphisme $P^\infty_{X/S} \simeq \mathcal{O}_{\hat P}$
prolongeant l'isomorphisme canonique d'ordre $2$. Ces isomorphismes sont en
bijection avec les relèvements à l'ordre $3$, et forment donc un torseur
canonique
\[
  \mathcal{C}_{X/S} \quad\text{sous}\quad
  \underline{\mathrm{Hom}}(\omega, \omega^3) \simeq \omega^2 :
\]
le \emph{torseur des connexions projectives}, qui est le
$\mathrm{Conn}_{X/S}$ de la station IX. De façon équivalente, ce sont les
connexions sur $\mathfrak{g} \simeq \mathrm{Sym}^2 E \otimes \omega^{-1}$ (gradué
par $\omega, \mathcal{O}, \omega^{-1}$) qui se restreignent sur $\omega$ en
l'opérateur universel $d^1_\omega : \omega \to P^1_{X/S}(\omega)$ — « on
retombe sur ses pattes ». La construction transcendante est donc celle d'une
\emph{section canonique} de $\mathcal{C}_{X/\mathbb{C}}$.

\textbf{Régularité à l'infini} (pages 134 à 137). Soit $\hat X$ la complétée
lisse de $X$, $D = \hat X - X$ (multiplicités $1$, $\nu = \mathrm{card}\,D$),
$\lambda = \omega^2_{\hat X}$. « Sauf erreur », la régularité à l'infini doit
signifier que la section n'a que des pôles simples le long de $D$, le groupe
des indéterminations des connexions projectives régulières étant
$H^0(\hat X, \omega^2(D))$.\footnote{Lu au sens strict, « régulière » permet
des pôles doubles : les connexions projectives à singularités régulières le
long de $D$ forment un torseur sous $\omega^2(2D)$. Le groupe
$H^0(\omega^2(D))$ est celui des connexions dont le terme dominant en chaque
point de $D$ est fixé — par exemple les connexions de type « pointe »,
comme la connexion uniformisante — et c'est lui que les comptes qui suivent
utilisent. On lit l'énoncé en ce sens.} La suite exacte
\[
  0 \to H^0(\hat X, \lambda) \to H^0(\hat X, \lambda(D)) \to
  H^0(D, \omega|_D) \to H^1(\hat X, \lambda)
\]
donne les dimensions :
\begin{enumerate}
  \item[(a)] si $g \geq 2$, $H^1(\hat X, \omega^2) = 0$, et l'espace des
  connexions projectives régulières est de dimension $(3g-3) + \nu$, avec
  $H^0(\hat X, \omega^2)$ cotangent à la variété modulaire $M_g$ et
  $H^0(X, \omega^2)^{\text{rég}}$ « (canoniquement ?) » cotangent à
  $M_{g,\nu}$ ;
  \item[(b)] si $g = 1$ et $\nu \geq 1$, le défaut d'injectivité passe par
  $H^0(\hat X, T_{\hat X})$, de dimension $1$, et la dimension est $\nu$ ;
  \item[(c)] si $g = 0$, $\dim H^0(\hat X, T) = 3$ et $\dim H^0(\hat X,
  T(-D)) = 2, 1, 0$ pour $\nu = 1, 2, \geq 3$ ; d'où
  $\dim H^0(X, \omega^2)^{\text{rég}} = \max(0, \nu - 3)$.
\end{enumerate}
\footnote{Dans le cas $g = 1$, il cherche une interprétation modulaire
(courbes elliptiques avec $\nu$ points « de somme nulle »), lue de façon très
incertaine. Dans le cas $g = 0$, la transcription relève une contradiction
apparente entre « $0 \leq \nu \leq 3$ » et le cas $\nu \geq 3$ : il n'y en a
pas, le conoyau étant de dimension $\nu$ exactement pour $\nu \leq 3$, et $3$
ensuite.}
Il n'y a donc qu'\emph{une seule} connexion projective régulière sur une
courbe de genre $0$ à au plus trois trous — celle qu'induit la droite
projective, qui est par conséquent la connexion uniformisante —, et pour
$\nu \geq 4$ un espace de dimension $\nu - 3$, la dimension de l'espace
modulaire $M_{0,\nu}$, « et sans doute un isomorphisme canonique »
$H^0(X, \omega^2)^{\text{rég}} \simeq \Omega^1_{M_{0,\nu}}(s)$.\footnote{C'est
bien l'espace cotangent, par la théorie des déformations des courbes
pointées. Ces deux pages répondent à la question de la page 114 pour la
droite moins trois points. La fin de la page 137 est en suspens.}

\pagerange{138}{139}
\subsection*{Connexions invariantes sur un torseur (pages 138 et 139)}

Deux pages très incertaines, sous l'hypothèse « inversible sur $S$ » d'un
entier qu'on ne lit pas : $H$ un schéma en groupes lisse de dimension
relative $1$ sur $S$, $X$ un torseur sous $H$. Les connexions projectives
invariantes par $H$ correspondent aux sections d'un sous-torseur de
$\mathfrak{g}_0 \otimes \omega_0$ sous $\omega_0^2$, où $\mathfrak{g}_0 \simeq
\omega_0 \oplus \mathcal{O}_S \oplus \omega_0^{-1}$ est l'algèbre de Lie descendue
sur $S$ ; parmi elles il y a une connexion canonique, caractérisée sur un
corps algébriquement clos par l'existence d'une section invariante de $P$
horizontale. Si $H$ est propre (schéma en courbes elliptiques), toutes les
connexions projectives sont invariantes.\footnote{Cette dernière affirmation
se vérifie : sur une courbe elliptique, $\omega$ est trivial et
$H^0(\omega^2)$ est formé de formes invariantes. Le reste de ces deux pages
— le rôle d'une formule « de Campbell--Hausdorff », le nombre « exactement
deux » de sections invariantes pour les autres connexions — est lu par
conjecture, mot à mot, et on ne s'en porte pas garant.}

\pagerange{141}{147}
\section*{XI. Connexions, stratifications, et équations de Picard--Fuchs (pages 141 à 160)}

La couverture brune qui ferme cette chemise (page 160) porte : « Modules
stratifiés et équations de Picard-Fuchs », avec « Connexions intégrables et »
biffé. Les sections sont numérotées par lui de 1) à 4).

\subsection*{1) Connexions sur les torseurs affines (pages 141 à 147)}

Soit $q : S \to T$ et $E$ un $\mathcal{O}_S$-module quasi-cohérent. On a
l'extension $0 \to E \otimes \Omega^1_{S/T} \to P^1_{S/T}(E) \to E \to 0$ et
son scindage $q^{-1}\mathcal{O}_T$-linéaire $d^1_E$, qui identifie additivement
$P^1_{S/T}(E)$ à $E \oplus E \otimes \Omega^1_{S/T}$, avec
$f(d^1_E x + \eta) = d^1_E(fx) + (f\eta - x \otimes df)$. Une \emph{connexion}
est un scindage $\mathcal{O}_S$-linéaire $\delta_E = d^1_E - d^0_E$ ; la condition
porte sur $d^0_E : E \to E \otimes \Omega^1_{S/T}$, qui doit vérifier
$d^0_E(fx) = f\,d^0_E(x) + x \otimes df$. Elle s'étend en un opérateur de degré
$1$ sur $E \otimes \Omega^\bullet_{S/T}$, de carré nul si et seulement si la
\emph{courbure} $d^1_E d^0_E$ est nulle ; on définit alors
$H^*_{\mathrm{DR}}(S,E) = \mathbb{H}^*(S, E \otimes \Omega^\bullet_{S/T})$, avec la
suite spectrale de la filtration bête, les flèches
$H^*\Gamma(S, E\otimes\Omega^\bullet) \to H^*_{\mathrm{DR}}(S,E)$ (isomorphismes
pour $S$ affine) et $H^i_{\mathrm{DR}}(S,E) \to H^i(S,E)$ ; $H^0_{\mathrm{DR}}$ est
le module des sections horizontales.

\textbf{$\mathbb{H}^1$ par les torseurs.} Pour un complexe de faisceaux
abéliens $K^\bullet$ sur un topos, $\mathbb{H}^1(K^\bullet)$ classe les couples
$(P, \xi)$ : $P$ un torseur sous $K^0$, $\xi$ une section de
$P \times^{K^0} K^1$ dont l'image dans $P \times^{K^0} K^2$ est la section
canonique. De façon équivalente, $\xi$ est une application
$\bar d^0 : P \to Z^1K$ telle que $\bar d^0(x + f) = \bar d^0(x) + d^0(f)$.
\footnote{La page note la loi de $K^1$ multiplicativement :
« $\bar d^0(xf) = \bar d^0(x)\,d^0(f)$ ». On écrit les groupes additivement.}
Appliqué au complexe de De~Rham :

\textbf{Proposition.} $H^1_{\mathrm{DR}}(S,E)$ classe les torseurs $P$ sous $E$
munis d'une application $d_P : P \to E \otimes \Omega^1_{S/T}$ telle que
$d_P(x + \xi) = d_P(x) + d^0_E(\xi)$ et $d^1_E d_P = 0$ — une \emph{connexion
intégrable} sur $P$, compatible à celle de $E$. Les automorphismes d'un tel
objet sont les translations par les sections horizontales de
$E$.\footnote{La page écrit « $f \in \mathcal{O}_S(\cdot)$ » dans la première
condition ; c'est une section $\xi$ de $E$ qu'il faut. Sans la condition
d'intégrabilité, les connexions sur un torseur donné existent localement et
forment un torseur sous $E \otimes \Omega^1$ (page 145).}

\pagerange{147}{153}
\subsection*{2) et 3) Stratifications, et comparaison (pages 147 à 153)}

Une stratification de $E$ relative à $T$ est un $\delta^\infty_E : E \to
P^\infty_{S/T}(E)$ satisfaisant les conditions de cocycle, d'où
$H^*_{\mathrm{strat}}(S,E)$, cohomologie du site stratifiant, et les flèches
$H^*_{\mathrm{strat}} \to H^*_{\mathrm{DR}}$. Un torseur stratifié sous $E$ est un
torseur ordinaire $P$ muni d'une stratification compatible, et
$P^\infty_{S/T}(E)$, vu comme pro-faisceau, a une stratification canonique.
D'où, de la suite exacte $0 \to E \xrightarrow{\delta} P^\infty_{S/T}(E)
\xrightarrow{p} P^\infty_{S/T}(E)/\delta E \to 0$ :

\textbf{Proposition} (page 150). Pour un torseur $P$ sous $E$, une
stratification compatible équivaut à une application
$\lambda : P \to (P^\infty_{S/T}(E)/\delta E)^{\natural}$ — à valeurs dans les
sections horizontales — telle que $\lambda(x + \xi) = \lambda(x) +
p\,d^\infty(\xi)$.\footnote{Le signe noté $\natural$ est, sur la page, « un petit
éclair » en exposant. L'encadré écrit « une stratification sur $E$ » : c'est
sur $P$.}

Par l'isomorphisme additif $P^\infty_{S/T}(E)^+ \simeq P^\infty_{S/T}(E)/\delta E$
(l'idéal d'augmentation), $\lambda$ devient $\lambda^+$ avec $\lambda^+(x +
\xi) = \lambda^+(x) + d^+(\xi)$, $d^+ = d^\infty - \delta^\infty$, et sa réduction
à l'ordre $1$ est la connexion associée.\footnote{La page 151 donne la
projection comme $\varphi \mapsto \varphi - \delta_E(\varphi)$ ; il faut
$\varphi - \delta_E(\varepsilon_E \varphi)$, $\varepsilon_E$ l'augmentation.
« Faire la vérification » est noté en marge.}

\textbf{Corollaire} (page 152). Si $T$ est de caractéristique nulle et $S$
lisse sur $T$, stratifications et connexions intégrables se correspondent, sur
$E$ comme sur les torseurs, et $H^*_{\mathrm{strat}} \simeq H^*_{\mathrm{DR}}$. En
particulier (le torseur trivial) la réduction induit un isomorphisme de
faisceaux $P^\infty_{S/T}(E)^{+\natural} \simeq Z^1(E \otimes
\Omega^\bullet_{S/T})$ : une $1$-forme fermée à coefficients dans $E$
\emph{est} un jet d'ordre infini intégrable.\footnote{La page dit
« $\Omega^1_{S/T}$ plat [p.ex.\ $S$ formellement lisse sur $T$] », et renvoie
à « [ ] » pour l'isomorphisme $H_{\mathrm{strat}} \simeq H_{\mathrm{DR}}$. La
platitude de $\Omega^1$ ne suffit pas à la correspondance
stratifications--connexions ; on énonce le cas lisse, celui de l'exemple, qui
est le théorème de son « Crystals and the de Rham cohomology of schemes »
(\emph{Dix exposés}, 1968).}

\pagerange{153}{159}
\subsection*{4) Équations de Picard--Fuchs (pages 153 à 159)}

Soit $E$ stratifié et $0 \to E' \xrightarrow{u} E \xrightarrow{v} E'' \to 0$
une suite exacte de modules quasi-cohérents, $E'$ \emph{non} stratifié. Soit
$P'$ un torseur sous $E'$ et $P = P' \times^{E'} E$. Une stratification de $P$
compatible à celle de $E$ équivaut à $\lambda' : P' \to (P^\infty(E)/\delta
E)^\natural$ avec $\lambda'(x' + \xi') = \lambda'(x') + p\,d^\infty(u\xi')$ ; en
poussant par $v$, on obtient une section canonique
\[
  \varphi(\lambda') \in \Gamma\bigl(S, (P^\infty_{S/T}(E'')/v\delta(E))^\natural\bigr),
\]
le \emph{jet de Picard--Fuchs} défini par $(P', \lambda')$ : un jet intégrable
de $E''$, défini seulement modulo l'image d'une section de $E'$.

Un opérateur différentiel $D : E'' \to F$ est un \emph{opérateur de
Picard--Fuchs} (relativement à la suite exacte) si l'application linéaire
$u_D : P^\infty(E'') \to F$ qu'il définit s'annule sur $v\delta(E)$ —
autrement dit, si $D$ tue les images dans $E''$ des sections horizontales de
$E$. On pose alors
\[
  \mu_D(P', \lambda') = \bar u_D\bigl(\varphi(\lambda')\bigr) \in \Gamma(S, F),
\]
la \emph{section de Manin} associée à $D$ et à $(P', \lambda')$ ; l'équation
universelle est à valeurs dans le pro-faisceau
$P^\infty(E'')/v\delta u(E')$.\footnote{Le nom est le sien. C'est la
construction de Yu.~Manin (1963) dans sa démonstration de la conjecture de
Mordell sur les corps de fonctions : appliquer aux intégrales abéliennes un
opérateur de Picard--Fuchs qui annule les périodes. On l'appelle aujourd'hui
« application de Manin ». Le transcripteur lit « $(P,\sigma)$ » là où il faut
$(P',\lambda')$ : on harmonise.}

Variante, si $E$ n'a qu'une connexion intégrable : les connexions intégrables
sur $P$ sont les $d_{P'} : P' \to E \otimes \Omega^1_{S/T}$ avec $d_{P'}(x' +
\xi') = d_{P'}(x') + d^0_E(u\xi')$ et $d^1_E d_{P'} = 0$ ; on prend des
opérateurs $\Delta : E \otimes \Omega^1_{S/T} \to F$ nuls sur $d^0_E(u(E'))$,
modulo ceux de la forme $\Delta' \circ d^1_E$, d'où une section
$\nu_\Delta(P', d_{P'}) \in \Gamma(S,F)$. Pour $D_0 = \Delta \circ d^0_E$, qui
passe au quotient en un opérateur de Picard--Fuchs $D$ sur $E''$, on a
$\nu_\Delta = \mu_D$ (« faire la vérification »).

\textbf{Lemme} (page 159). En caractéristique nulle, $S$ lisse sur $T$, les
opérateurs différentiels $E \to F$ dont l'application linéaire associée tue
$\delta(E)$ sont exactement les $\Delta \circ d^0_E$, $\Delta : E \otimes
\Omega^1 \to F$ différentiel, à un $\Delta' \circ d^1_E$ près. Il en résulte
que, sur $S$ affine, les $\mu_D$ s'expriment par les $\nu_\Delta$. Le lemme
est un cas particulier, dit-il, d'un fait général sur les opérateurs
différentiels, qu'il formule à la station XX.

\pagerange{161}{166}
\section*{XII. Le complexe de De~Rham filtré, et la connexion de Gauss--Manin (pages 161 à 166)}

Le titre est le sien : « Filtration du complexe de De Rham d'un schéma
relatif lisse, et connexion de Gauss-Manin (d'après Katz-Oda) ».\footnote{N.~Katz
et T.~Oda, « On the differentiation of de Rham cohomology classes with
respect to parameters », \emph{J. Math. Kyoto Univ.} 8 (1968).}

Soit $X \xrightarrow{f} S \xrightarrow{g} T$, $f$ lisse. La suite exacte
$0 \to f^*\Omega^1_{S/T} \to \Omega^1_{X/T} \to \Omega^1_{X/S} \to 0$ filtre
$\Omega^\bullet_{X/T}$, avec
\[
  \mathrm{Gr}^p(\Omega^\bullet_{X/T}) \simeq f^*\Omega^p_{S/T} \otimes
  \Omega^\bullet_{X/S}[-p] ,
\]
et la filtration est respectée par la différentielle. Par la formule de
projection (« dans un cadre pas très catholique — il faudra faire une petite
vérification »), pour $f$ quasi-compact et quasi-séparé,
$\mathrm{Gr}^p\,\mathbf{R}f_*\Omega^\bullet_{X/T} \simeq \Omega^p_{S/T}
\otimes^{\mathbf{L}} \mathbf{R}f_*\Omega^\bullet_{X/S}[-p]$. La différentielle
$d_1$ de la suite spectrale de la filtration donne, pour $p = 0$,
\[
  \nabla : \mathbf{R}f_*\Omega^\bullet_{X/S} \longrightarrow \Omega^1_{S/T}
  \otimes \mathbf{R}f_*\Omega^\bullet_{X/S} ,
\]
la \emph{connexion de Gauss--Manin}, intégrable puisque $d_1 d_1 = 0$ ; si
$\Omega^1_{S/T}$ est plat, on obtient une connexion intégrable sur chaque
$H^q_{\mathrm{DR}}(X/S) = \mathbf{R}^q f_*\Omega^\bullet_{X/S}$. « Il semble bien
que c'est la même que celle que j'ai définie » — dans un texte dont la
référence ne se lit pas.\footnote{Sans doute la définition cristalline,
par l'indépendance de la cohomologie de De~Rham vis-à-vis du relèvement,
que la station XXI reconstruit. Une marge avertit qu'il faut une
compatibilité de $d_1$ à la différentielle de $\Omega^\bullet_{S/T}$ pour
conclure à l'intégrabilité.}

La filtration donne la \emph{suite spectrale de Leray filtrée} (« lecture
incertaine du dernier mot »)
\[
  H^{p+q}_{\mathrm{DR}}(X/T) \Longleftarrow E_2^{p,q} =
  H^p_{\mathrm{DR}}\bigl(S/T, H^q_{\mathrm{DR}}(X/S)\bigr)
  \qquad (\Omega^1_{S/T} \text{ plat}),
\]
dont le terme $E_2$ se calcule avec la connexion de Gauss--Manin.

\textbf{Application.} Une classe $\xi \in H^n_{\mathrm{DR}}(X/T)$ a une image
horizontale dans $H^0(S, H^n_{\mathrm{DR}}(X/S))$ ; si cette image est nulle,
$\xi$ définit un élément de $E_\infty^{1,n-1} \subset E_2^{1,n-1} =
H^1_{\mathrm{DR}}(S/T, H^{n-1}_{\mathrm{DR}}(X/S))$, c'est-à-dire (station XI) une
classe de torseurs sous $H^{n-1}_{\mathrm{DR}}(X/S)$ munis d'une connexion
intégrable compatible à celle de Gauss--Manin.\footnote{La formule de
$E_2^{0,n}$ que la page porte en marge est surchargée et en partie biffée ;
on la remplace par $H^0_{\mathrm{DR}}(S/T, H^n_{\mathrm{DR}}(X/S))$, et l'on lit
$\Omega^1_{S/T}$ là où elle écrit $\Omega^1_{X/S}$. C'est la construction de
la classe d'extension (ou « fonction normale ») d'un cycle homologue à zéro
dans les fibres ; la page ne lui donne pas de nom.} Il note qu'en
caractéristique nulle, sans hypothèse de platitude, la donnée équivaut à une
stratification du torseur (station XI).

\emph{Variantes} (page 166) : des coefficients $E$ à connexion intégrable sur
$X$ — il suffit, semble-t-il, que $E$ soit plat sur $X$ ou $\Omega^1_{S/T}$
plat, d'où $H^*_{\mathrm{DR}}(X/T, E) \Leftarrow H^p_{\mathrm{DR}}(S/T,
H^q_{\mathrm{DR}}(X/S, E))$ ; et la relativisation sur une base au-dessus de $T$.

\pagerange{168}{172}
\section*{XIII. L'opération de Cartier (pages 167 à 186)}

La chemise porte « Opération de Cartier » (page 167). « Tous les anneaux sont
des $\mathbb{F}_p$-algèbres. »

\subsection*{① à ③. La cohomologie de De~Rham d'une algèbre à $p$-base, et l'isomorphisme de Cartier (pages 168 à 172)}

Soit $B$ une $A$-algèbre et $\Lambda = A[B^p] \subset B$. Alors
$\Omega^\bullet_{B/A} = \Omega^\bullet_{B/\Lambda}$. Si $B = \Lambda[X]/(X^p - a)$
(une $p$-base à un élément $x$), le complexe $\Lambda$-linéaire $1, x, \ldots,
x^{p-1} \to dx, \ldots, x^{p-1}dx$ envoie $x^k$ sur $kx^{k-1}dx$, et
$H^*_{\mathrm{DR}}(B/A)$ a pour base $1$ et la classe de $x^{p-1}dx$. Pour une
$p$-base $(x_i)_{i \in I}$, Künneth donne $H^*_{\mathrm{DR}}(B/A) \simeq
\bigotimes_{i,\Lambda} H^*_{\mathrm{DR}}(B_i/\Lambda)$.

Il existe un unique homomorphisme d'algèbres graduées $\psi : \Omega^*_{B/A}
\to Z^*_{B/A}$ avec $\psi(x) = x^p$ et $\psi(dx) = x^{p-1}dx$ (il suffit que
$(xy)^{p-1}d(xy) = x^p y^{p-1}dy + y^p x^{p-1}dx$). Notons $\widetilde{A} = A$
vu comme $A$-algèbre par Frobenius et $\widetilde{B} = B \otimes_A \widetilde A$,
dont l'image dans $B$ est $\Lambda$. Par passage aux scalaires on obtient
$\varphi : \Omega^*_{\widetilde B/\widetilde A} \to H^*_{\mathrm{DR}}(B/A)$.

\textbf{Théorème 1.} Si $B/A$ admet une $p$-base, $\varphi$ induit un
isomorphisme $\Omega^*_{\widetilde B/\widetilde A} \otimes_{\widetilde B} \Lambda
\xrightarrow{\sim} H^*_{\mathrm{DR}}(B/A)$ ; c'est un isomorphisme si $\widetilde B
\to B$ est injectif. (Par Künneth on se ramène à une $p$-base à un élément,
où $\varphi$ envoie la base $1, d\tilde x$ sur la base $1, [x^{p-1}dx]$.)

\textbf{Corollaire 1.} $Z^*(B/A) = B^*(B/A) \oplus \widetilde\psi\bigl(\Omega^*_{\widetilde
B/\widetilde A} \otimes \Lambda\bigr)$, où $B^*(B/A)$ désigne les formes exactes.

L'isomorphisme inverse est l'\emph{opération de Cartier}
$\widetilde C : H^*_{\mathrm{DR}}(B/A) \to \Omega^*_{\widetilde B/\widetilde A}
\otimes \Lambda$, et l'on note encore $\widetilde C : Z^*(B/A) \to
\Omega^*_{\widetilde B/\widetilde A}\otimes\Lambda$ le composé, surjectif de noyau
les formes exactes. Il est caractérisé (corollaire 2) par : $\widetilde C$ est
nul sur les formes exactes, est un homomorphisme de $\Lambda$-algèbres, et
$\widetilde C(x^{p-1}dx) = d\tilde x \otimes 1$. Pour $x$ inversible,
$\widetilde C(dx/x) = d\tilde x/\tilde x \otimes 1$ (corollaire 3, car $dx/x =
x^{-p}\cdot x^{p-1}dx$).\footnote{Dans le diagramme de la page 169, l'arc est
étiqueté $F_A$ là où l'on attendrait $F_B$ ; un coin du feuillet manque, et
la fin de quelques lignes est perdue.} Si $A$ est parfait et $B$ réduit, on
retrouve la forme classique : $C : Z^*(B) \to \Omega^*_B$,
$C(f^p\omega) = fC(\omega)$, $C(x^{p-1}dx) = dx$, $C(dx/x) = dx/x$ — « le cas
envisagé dans le texte initial de Cartier ».

\pagerange{172}{180}
\subsection*{④ à ⑥. Globalisation ; différentielles logarithmiques ; formes de degré maximum (pages 172 à 180)}

Pour $X \to S$ en caractéristique $p$, avec le Frobenius relatif
$F_{X/S} : X \to X' = X^{(p/S)}$ (radiciel, entier, surjectif), on a
$\psi : \Omega^*_{X'/S} \to F_{X/S*}\mathcal{H}^*(\Omega^\bullet_{X/S})$, un
isomorphisme quand $F_{X/S*}\mathcal{O}_X$ a localement une $p$-base, d'inverse
$\widetilde C$ ; la construction est fonctorielle.

\textbf{Connexions sur un module inversible} (pages 173 à 175). Soit $L$
inversible muni d'une connexion $\Delta_0$ à courbure nulle et compatible aux
puissances $p$-ièmes. Les connexions sur $L$ s'écrivent $\Delta = \Delta_0 +
\omega$ ; elles sont à courbure nulle si et seulement si $d\omega = 0$, et alors
compatibles aux puissances $p$-ièmes si et seulement si
\[
  \widetilde C\omega = \tilde\omega \quad (\text{l'image de } \omega \text{ par le
  changement de base de Frobenius}).
\]
En coordonnées, avec une $p$-base $(x_i)$, $D_i = \partial/\partial x_i$ et
$\omega = \sum c_i\,dx_i$, la formule de Jacobson donne $(\Delta^0_i + c_i)^p =
c_i^p + D_i^{p-1}c_i$, et la condition est $c_i^p + D_i^{p-1}c_i = 0$ ; or
$\widetilde C\omega = \sum -D_i^{p-1}c_i\,(d\tilde x_i \otimes 1)$ et
$\tilde\omega = \sum c_i^p\,(d\tilde x_i \otimes 1)$.\footnote{Il renvoie au
Séminaire Chevalley 1958/59, exposé 6, pour la formule de Jacobson. La
version en rang supérieur, $d\omega + [\omega,\omega] = 0$ pour la courbure,
est notée entre crochets. La fin de la page 173 est une superposition
d'insertions dont l'ordre reste incertain.}

Comme une connexion à courbure nulle et compatible aux puissances $p$-ièmes
équivaut à une donnée de descente le long de $X \to X'$, on obtient :

\textbf{Théorème.} Une $1$-forme $\omega$ est localement une différentielle
logarithmique $df/f$, $f \in \mathcal{O}_X^*$, si et seulement si $d\omega = 0$
et $\widetilde C\omega = \tilde\omega$. Le noyau de $f \mapsto df/f$ est
$\Lambda^*$, soit $\mathcal{O}_X^{*p}$ si $S$ est parfait.

\textbf{Le cas lisse} (page 176). Si $X$ est lisse sur $S$, $F_{X/S}$ est fini
localement libre, $X$ a localement une $p$-base, et
$\widetilde C : F_{X/S*}(Z\Omega^*_{X/S}) \to \Omega^*_{X'/S}$.

\textbf{Théorème 3.} En degré maximum $n$ (la dimension relative), $\widetilde
C : F_{X/S*}\Omega^n_{X/S} \to \Omega^n_{X'/S}$ est le morphisme trace (de
Gysin) de $F_{X/S}$ pour la dualité. De façon équivalente, pour un
sous-schéma fini défini par $f_1, \ldots, f_n$,
\[
  \operatorname{res}\Bigl(\frac{\widetilde C\omega}{\tilde f_1 \cdots \tilde f_n}\Bigr)
  = \operatorname{res}\Bigl(\frac{\omega}{f_1^p \cdots f_n^p}\Bigr),
\]
et, en dimension $1$ sur une base parfaite, $(\operatorname{res} C\omega)^p =
\operatorname{res}\omega$.\footnote{C'est la formule de J.~Tate (« Residues of
differentials on curves », \emph{Ann. Sci. ÉNS} 1, 1968). Le théorème est
énoncé, non démontré.}

\emph{Remarques.} ① L'image inverse $F^* : \Omega^i_{X'} \to F_*\Omega^i_X$ est
nulle pour $i \geq 1$, de sorte que les morphismes de type « trace » qu'on
définirait par dualité en degré $i \neq n$ ne sont pas l'opération de
Cartier sur les formes fermées. ② On peut définir $C$ sur le module dualisant
relatif sans hypothèse de lissité, lorsque sa formation commute au
changement de base par $F_S$ (par exemple $S$ parfait).

\textbf{Dualité} (pages 179 et 180). Pour $X$ propre et lisse sur $S$ parfait et
$E$ localement libre, le transposé (dualité de Serre) du morphisme
« puissance $p$-ième » $\mathbf{R}f_*(\widetilde E) \to \mathbf{R}f_*(E^{(p)})$ est
induit par la trace $F_{X/S*}(\check E^{(p)} \otimes \Omega^n) \to \check{\widetilde E}
\otimes \Omega^n_{X'}$, c'est-à-dire par l'opération de Cartier. Pour $E =
\mathcal{O}_X$ : le Frobenius $R^if_*(\mathcal{O}_X)^{(p)} \to R^if_*\mathcal{O}_X$
est le transposé de $\widetilde C$ sur $R^{n-i}f_*\Omega^n_{X/S}$.\footnote{C'est
la dualité entre matrice de Hasse--Witt et matrice de Cartier--Manin. La page
écrit « $f_*(\Omega^n_{X/S})^{\otimes p}$ » où il faut le tordu
$(\,)^{(p)}$, et met plusieurs isomorphismes « sous réserve » de commutation
au changement de base par $F_S$.} Pour $i = 1$, $R^1f_*\mathcal{O}_{X'} \to
R^1f_*\mathcal{O}_X$ est l'application tangente de $F^* : \mathrm{Pic}_{X'/S}
\to \mathrm{Pic}_{X/S}$, et l'on « admet » que celle-ci est le Verschiebung
$V : \mathrm{Pic}^{(p/S)}_{X/S} \to \mathrm{Pic}_{X/S}$ (avec $FV = VF =
p$).\footnote{Pour un schéma abélien, c'est le fait que le dual du Frobenius
relatif est le Verschiebung du dual. La marge propose d'en tirer des
conditions sur le $p$-rang.}

\pagerange{181}{186}
\subsection*{⑦. Courbes : un accouplement remarquable (pages 181 à 186)}

Supposons $\Omega^1_{X/S}$ inversible, c'est-à-dire que $X$ a localement une
$p$-base à un élément sur $S$, et $Q = F_{X/S*}\mathcal{O}_X/\Lambda$, localement
libre de rang $p - 1$. On pose
\[
  \varphi(f,g) = \widetilde C(f\,dg) \in \Omega^1_{X'/S} \otimes \Lambda .
\]
Cette forme est $\Lambda$-bilinéaire, nulle si $f$ ou $g$ est dans $\Lambda$,
et antisymétrique ($\varphi(f,g) + \varphi(g,f) = \widetilde C(d(fg)) = 0$) :
alternée si $p \neq 2$, symétrique si $p = 2$.

\textbf{Proposition.} L'accouplement induit $Q \times Q \to \Omega^1_{X'/S}
\otimes \Lambda$ est parfait. En effet, pour une $p$-base $x$, la base
$x^i$ ($1 \leq i \leq p-1$) de $Q$ donne $\widetilde C(x^i\,dx^j) =
\widetilde C(j\,x^{i+j-1}dx)$, égal à $j\,d\tilde x$ si $i + j = p$ et nul
sinon.\footnote{La page écrit « si $i + j = p - 1$ » ; la condition est
$i + j - 1 = p - 1$, soit $i + j = p$, ce que confirme la matrice
antidiagonale qui suit.} La matrice est antidiagonale de coefficients $1, 2,
\ldots, p-1$, de déterminant $\pm (p-1)! = \mp 1$ par le théorème de Wilson.

\textbf{Corollaire.} $Q \simeq \check Q \otimes (\Omega^1_{X'/S} \otimes \Lambda)$,
d'où $(\det Q)^{\otimes 2} \simeq (\Omega^1_{X'/S})^{\otimes(p-1)} \otimes
\Lambda$.

\emph{Cas $p = 2$} : $\Omega^1_{X'/S} \otimes \Lambda$ est un carré, celui de
$\det Q$ — en particulier, sur un corps et pour $X$ réduit,
$\Omega^1_{X/S}$ est lui-même un carré.\footnote{C'est la « thêta-caractéristique
canonique » d'une courbe en caractéristique $2$. Le passage est surchargé ; il
y évoque un contre-exemple sur corps imparfait (« Serre-Tate-Fröhlich »,
lecture incertaine) et une courbe rationnelle sans point rationnel.}

\emph{Cas $p \neq 2$} : on aimerait relever $(**)$ en un isomorphisme canonique
$\psi : \det Q \simeq (\Omega^1_{X'/S})^{\otimes\frac{p-1}{2}} \otimes \Lambda$,
défini par $\psi(\dot x \wedge \dot x^2 \wedge \cdots \wedge \dot x^{p-1}) =
d\tilde x^{\otimes\frac{p-1}{2}}$, et il faut vérifier que ce choix est
compatible à $(**)$ et indépendant de la $p$-base ; pour l'indépendance, il
esquisse un argument de connexité des fibres du schéma des $p$-bases
relatives.\footnote{La compatibilité demande que le déterminant $\Delta$ de la
matrice vaille $1$. La page le calcule comme $(-1)^{p-2}(p-1)!$ et conclut
$\Delta = 1$. Mais la signature de la permutation antidiagonale de $p - 1$
lettres est $(-1)^{(p-1)(p-2)/2}$, d'où, avec les identifications usuelles,
$\Delta = (-1)^{(p+1)/2}$ : $1$ pour $p \equiv 3 \pmod 4$, $-1$ pour $p \equiv
1 \pmod 4$. Dans ce second cas $-1$ est un carré dans $\mathbb{F}_p$ et un
relèvement existe encore, mais la normalisation $(***)$ doit être changée
d'une racine carrée de $-1$, ce qui réintroduit l'ambiguïté de signe que
$\psi$ devait lever. Le signe dépend aussi de la convention qui identifie
$\det\check Q$ à $(\det Q)^{-1}$ ; nous ne tranchons pas. Que $\det
F_*\mathcal{O}_X$ soit, en degré, $\omega^{(p-1)/2}$ se vérifie par
Riemann--Roch.}

\textbf{Le diagramme autodual} (page 185). Pour $X \to S$ propre et lisse de
dimension relative $1$, $H^1_{\mathrm{DR}}(X/S)$ porte deux filtrations : la
filtration de Hodge (la colonne), $0 \to f_*\Omega^1 \to H^1_{\mathrm{DR}} \to
R^1f_*\mathcal{O}_X \to 0$, et une ligne
\[
  0 \to R^1f_*(\mathcal{O}_X)^{(p)} \xrightarrow{\ \alpha\ } H^1_{\mathrm{DR}}
  \xrightarrow{\ \beta\ } f_*(\Omega^1_{X/S})^{(p)} \to 0 ,
\]
$\alpha$ déduit du Frobenius de $H^1$, $\beta$ son transposé, les flèches
composées étant le Frobenius $F$ de $R^1f_*\mathcal{O}$ et l'opération de
Cartier $C$, transposées l'une de l'autre par la dualité de Poincaré. On a
toujours $\operatorname{Ker}\alpha = (\operatorname{Im}\beta)^\perp$ et
$\operatorname{Im}\alpha = (\operatorname{Ker}\beta)^\perp$.\footnote{La ligne est
aujourd'hui la \emph{filtration conjuguée} (N.~Katz), dont l'exactitude, pour
une courbe, vient de ce que la suite spectrale conjuguée dégénère par raison
de degré. La page, après avoir écrit « ligne et colonne exactes », barre un
encadré qui cherchait des conditions équivalentes à cette exactitude.} La
page 186, scannée à l'envers et barrée, porte des diagrammes d'une autre
nature (immersions et voisinages infinitésimaux) ; on ne la restitue pas.

\pagerange{188}{192}
\section*{XIV. Connexions et stratifications intégrables : problèmes divers (pages 187 à 197)}

Le titre est celui de la couverture (page 187). Une marge note qu'« il serait
plus naturel » de travailler avec les PD-connexions, « pareil dans le cas
lisse ».

\textbf{1) Connexions triviales.} Un module à connexion $(E, \nabla)$ sur $X/S$
est \emph{trivial} s'il est isomorphe à $f^*F$ avec sa connexion canonique ;
si $X \to S$ est lisse, surjectif, à fibres géométriquement connexes et $S$ de
caractéristique nulle, $F \mapsto f^*F$ est pleinement fidèle. Il est
\emph{intégrable} s'il est trivial localement pour la topologie étale
(« Zar.-intégrable » pour celle de Zariski).

\textbf{2) Caractéristique $p$.} Si $X/S$ a localement une $p$-base, sont
équivalents : $E$ est Zariski-intégrable ; $E$ est intégrable ; $E$ est à
courbure et $p$-courbure nulles ; $E$ provient d'un module sur $X'$ (ou sur
$(X', \Lambda)$). Et le foncteur $F \mapsto F^*_{X/S}F$ est une équivalence des
modules quasi-cohérents sur $(X', \Lambda)$ vers les modules à connexion
intégrable.\footnote{C'est le théorème de descente de Cartier, sous la forme
que lui donne N.~Katz (« Nilpotent connections and the monodromy
theorem », \emph{Publ. Math. IHÉS} 39, 1970, §5).}

\textbf{3) Stratifications et groupe fondamental.} Un torseur sous un groupe
étale $G$ et une représentation de $G$ produisent, par torsion, un module à
stratification intégrable. Sur un corps $k$, avec un point rationnel $x$, la
catégorie des modules à stratification intégrable « doit être équivalente »
à celle des représentations de dimension finie d'un pro-groupe
$\underline{\Pi}_1(X,x)$ ; et sur une base générale, pour un ensemble
$\mathbb{L}$ de nombres premiers, à celle des représentations d'un
$\underline{\Pi}_1^{\mathbb{L}}(X,e)$ dans des modules cohérents.\footnote{Le
groupe qui convient n'est pas le groupe fondamental étale : en
caractéristique $p$, on parle aujourd'hui du schéma en groupes fondamental
\emph{stratifié}, défini par le formalisme tannakien (N.~dos~Santos,
\emph{J. Algebra}, 2007) ; voir aussi D.~Gieseker (1975), qui conjectura que
les fibrés stratifiés sur une variété projective lisse simplement connexe
sont triviaux, et H.~Esnault et V.~Mehta (\emph{Invent. Math.}, 2010), qui
l'ont démontré. Les pages, très corrigées, sont lues en partie.}

\textbf{4) Le rang un.} Pour $X$ de type fini sur un corps, géométriquement
unibranche, un module inversible stratifié est intégrable si et seulement
si une puissance $L^{\otimes n}$, $n$ premier à $p$, est triviale comme module
stratifié. En caractéristique nulle, les stratifications de $\mathcal{O}_X$ sont
les formes fermées $\omega$ ; la triviale correspond aux différentielles
logarithmiques ; l'intégrabilité, à ce qu'un multiple $n\omega$ en soit une.
Comme le groupe des différentielles logarithmiques est de type fini
($\simeq \mathbb{Z}^r$, par le théorème des unités de Rosenlicht), il y a des
connexions intégrables non triviales sur $\mathcal{O}_X$, et leurs classes
forment un groupe $\simeq (\mathbb{Q}/\mathbb{Z})^r$ ; si $\Gamma(X,
\mathcal{O}_X^*)$ est réduit aux constantes, il n'y en a pas d'autre que la
triviale.

\pagerange{193}{197}
\subsection*{Questions (pages 193 à 197)}

Soit $S$ intègre, de type fini et dominant sur $\mathrm{Spec}\,\mathbb{Z}$, de point
générique $\eta$, $X$ lisse de type fini sur $S$, et $E$ localement libre sur
$X$ avec une connexion à courbure nulle relative à $S$. La connexion de
$E_\eta$ est dite \emph{intégrable} si $E_\eta$ devient localement trivial sur
un revêtement fini étale de $X_\eta$ — sur $\mathbb{C}$, si sa monodromie est
finie. Si elle l'est, il existe un voisinage $U$ de $\eta$ tel que, pour
tout $s \in U$, la $p$-courbure de $E_s$ sur $X_s/k(s)$ soit nulle (il suffit
de le demander aux points fermés).

\emph{Question} : « la réciproque est-elle vraie ? »\footnote{C'est la
conjecture de Grothendieck--Katz sur la $p$-courbure. N.~Katz l'a démontrée
pour les connexions de Gauss--Manin (\emph{Invent. Math.} 18, 1972) ; le cas
où la monodromie est résoluble, qui contient le rang un, a été établi par
D.~et G.~Chudnovsky pour les courbes, puis par J.-B.~Bost (2001) et Y.~André
(2004). Le cas général est ouvert.}

\emph{Le rang un} se décompose en deux questions : a) si les connexions des
$E_s$ sont intégrables pour presque tout $s$, l'image de $E_\eta$ dans
$\mathrm{Pic}(X_\eta)$ est-elle de torsion ? b) si $\omega$ est une $1$-forme
fermée sur $X/S$ dont les réductions $\omega_s$ sont localement des
différentielles logarithmiques — par Cartier, $\widetilde C\omega_s =
\tilde\omega_s$ —, un multiple entier de $\omega_\eta$ est-il une
différentielle logarithmique ?\footnote{Deux paragraphes portent la lettre
b) ; le premier, qui réduit au cas $E \simeq \mathcal{O}_X$, est corrigé à la
main.}

Il les généralise aux torseurs sous un schéma en groupes $G$ (pages 195 à
197 ; la page 196 précède la page 195 dans l'ordre de la rédaction). Pour un
schéma abélien, avec l'extension vectorielle universelle $E$ du dual : c)
une section $\varphi$ de $E$ telle que $\varphi(s) \in pE_s$ pour tout point
fermé $s$ est-elle d'ordre fini ? d) pour un schéma en groupes de type fini
et une section $\omega$ qui tombe, en chaque point fermé, dans un sous-groupe
connexe d'un type donné (unipotent, affine, de rang unipotent nul), en
est-il de même au point générique ? Les cas « affine » et « unipotent »
seraient triviaux, par le théorème de Chevalley. e) si une section $\omega$
d'un schéma abélien se relève, en chaque point fermé, par le Verschiebung
$A_s^{(p)} \to A_s$, est-elle d'ordre fini ?\footnote{Ces questions sont lues
avec beaucoup d'incertitudes aux pages 195 et 196 ; on n'en garde que la
forme. Nous ne disons rien de leur état actuel.}

La page 197 se termine par « Fin 1ère partie ».

\pagerange{198}{203}
\section*{XV. Itérés de Cartier et valeurs propres de Frobenius (pages 198 à 203)}

Sous les conditions où $\widetilde C : f_*\Omega^1_{X/S} \to f_*\Omega^1_{X/S}{}^{(p)}$
est défini et $\mathcal{O}_S$-linéaire (par exemple $f_*\Omega^1 = f_*Z^1$),
notons $\underline{\omega} = f_*\Omega^1_{X/S}$ et formons
$F_*^{(\nu)} = \widetilde C^{(p^{\nu-1})} \cdots \widetilde C^{(p)}\widetilde C :
\underline\omega \to \underline\omega^{(p^\nu)}$. Si $\lambda \mapsto \lambda^{p^\nu}$
est l'identité sur $\mathcal{O}_S$ (par exemple $S$ sur $\mathbb{F}_{p^\nu}$), c'est
un endomorphisme $\mathcal{O}_S$-linéaire de $\underline\omega$, et pour une forme
logarithmique ($\widetilde C\omega = \omega^{(p)}$) on a $F_*^{(\nu)}\omega =
\omega$. En dimension relative $1$, $F_*^{(\nu)}$ est la trace du Frobenius
itéré $F^\nu_X$, qui est alors un $S$-endomorphisme de $X$ ; c'est le transposé
de l'action de $F^\nu_X$ sur la cohomologie.

Sur $S = \mathrm{Spec}\,\mathbb{F}_{p^\nu}$, $X$ propre, l'opération induite sur
$H^1_{\mathrm{DR}}/\underline\omega \simeq R^1f_*\mathcal{O}$ est nulle, et il en
conclut que les valeurs propres non nulles de $F_*^{(\nu)}$ sur
$\underline\omega$ sont, modulo $p$, celles des valeurs propres de Frobenius
($\ell$-adiques ou cristallines) qui sont des unités $p$-adiques —
« (Katz) ». D'où :

\textbf{Énoncé} (page 203). Il existe $\omega \neq 0$ avec $\widetilde C\omega =
\tilde\omega$ si et seulement si une valeur propre de Frobenius est congrue à
$1$ modulo $p$ ; la dimension de l'espace des formes logarithmiques est le
nombre de ces valeurs propres, et c'est aussi $\dim_{\mathbb{F}_p}
{}_p\mathbf{A}(k)$, $\mathbf{A}$ la jacobienne.\footnote{La première
équivalence est juste, la congruence s'entendant modulo un idéal premier
au-dessus de $p$ : ${}_p\mathbf{A}(k) \neq 0$ si et seulement si $p$ divise
$\#\mathbf{A}(k) = \prod(1 - \alpha_i)$. L'égalité des dimensions demande plus :
le nombre de valeurs propres congrues à $1$, comptées avec multiplicité, est
la dimension d'un sous-espace caractéristique, qui majore celle de l'espace
des points fixes ; il y a égalité si Frobenius agit de façon semi-simple sur
cette partie. Sur la clôture algébrique, l'espace des formes logarithmiques
a pour dimension le $p$-rang de la jacobienne. Les pages 202 et 203 sont
lues avec de nombreuses incertitudes.} « Fin de parenthèse ».

\pagerange{205}{212}
\section*{XVI. Foncteurs de Picard relatifs avec connexion (pages 204 à 212)}

Le titre est celui de la chemise (page 204). Soit $f : X \to S$ et
$f^{\mathrm{fl}}$ le morphisme des sites fppf. On définit, comme faisceaux
associés :
\begin{enumerate}
  \item[①] $P^c_{X/S}$ : les modules inversibles munis d'une connexion
  relative. On a la suite exacte
  \[
    0 \to f_*\Omega^1_{X/S}/d\log f_*\mathcal{O}_X^* \to P^c_{X/S} \to
    \underline{\mathrm{Pic}}_{X/S} \to R^1f_*\Omega^1_{X/S} ,
  \]
  la dernière flèche étant la classe d'Atiyah. Si $f$ est propre et
  $\underline{\mathrm{Pic}}^0_{X/S} = B$ un schéma abélien, tout morphisme de $B$
  vers $R^1f_*\Omega^1$ est nul, d'où $0 \to f_*\Omega^1_{X/S} \to (P^c)^0
  \to B \to 0$ ; pour une courbe propre et lisse, ou un torseur sous un schéma
  abélien, c'est l'\emph{extension vectorielle universelle} de $B$.
  \footnote{B.~Mazur et W.~Messing, LNM 370 (1974), où l'on note
  $\mathrm{Pic}^\natural$ ce foncteur. La page vérifie l'identification par
  $\mathrm{Lie}(B^*)^\vee \simeq f_*\Omega^1$, l'albanaise $A = B^*$ et
  $X \to \mathrm{Alb}^1$ ; le crochet ne se ferme pas.}
  \item[②] $P^{ci}$ : les connexions à courbure nulle ;
  $0 \to P^{ci} \to P^c \to f_*Z^2_{X/S}$, et dans le cas lisse
  $0 \to f_*Z^1/d\log \to P^{ci} \to \underline{\mathrm{Pic}} \to
  \mathbf{R}^2f_*(\Omega^{\geq 1}_{X/S})$. Si $f_*\Omega^1 = f_*Z^1$ (toutes les
  formes globales sont fermées), $(P^c)^0 = (P^{ci})^0$.
  \item[③] Si $p\cdot 1_S = 0 $ et $X/S$ a localement une $p$-base,
  $P^{cip} \subset P^{ci}$ : les connexions de $p$-courbure nulle. On a
  $0 \to P^{cip} \to P^{ci} \xrightarrow{K} f_*(\Omega^1)^{\otimes p}$, $K$ la
  $p$-courbure, et pour la connexion $d + \omega$ sur $\mathcal{O}$,
  \[
    K_0(\omega) = \omega^{\otimes p} - F^*_{X/S}(\widetilde C\omega)
    \qquad (d\omega = 0),
  \]
  additive, somme d'une partie $p$-linéaire et d'une partie linéaire. Par la
  descente de Cartier, $P^{cip}$ est le faisceau associé à $S' \mapsto
  \mathrm{Pic}((X_{S'})^{(p)})$ ; si $X$ est lisse, $P^{cip} \simeq
  \underline{\mathrm{Pic}}_{X'/S} \simeq \underline{\mathrm{Pic}}^{(p)}_{X/S}$, et
  $P^{cip} \to \underline{\mathrm{Pic}}$ est le Verschiebung. D'où
  \[
    \operatorname{Ker}(V : \underline{\mathrm{Pic}}^{(p)}_{X/S} \to
    \underline{\mathrm{Pic}}_{X/S}) \simeq
    \{\omega : d\omega = 0,\ C\omega = \tilde\omega\}/d\log f_*\mathcal{O}^* .
  \]
  \item[④] Si $p$ est seulement nilpotent sur $S$, $P^{cin}$ : les connexions
  à courbure nulle \emph{nilpotentes} ; sur les bases réduites ce sont les
  mêmes que $P^{cip}$, et $P^{cin}$ apparaît comme le complété formel de
  $P^{cip} \simeq \underline{\mathrm{Pic}}^{(p)}$ dans $P^{ci}$.
\end{enumerate}
\emph{Question} (page 208) : pour un module $E$ à connexion à courbure nulle
en caractéristique $p$, peut-on définir un tenseur de $p$-courbure
$K \in \Gamma(X, \underline{\mathrm{End}}(E) \otimes (\Omega^1)^{\otimes p})$
?\footnote{La réponse est oui (N.~Katz, 1970) : $\psi(D) = \nabla(D)^p -
\nabla(D^{[p]})$ est $p$-linéaire en $D$, et définit une section de
$\underline{\mathrm{End}}(E) \otimes F^*_{X/S}\Omega^1_{X'/S}$, ce qui ne coïncide avec
$(\Omega^1)^{\otimes p}$ qu'en dimension relative $1$.}

Pour un schéma abélien $B$, l'extension $P^{cip,0}$ de $B$ par
$\alpha_B = \operatorname{Ker}(V_B)$ serait l'extension universelle de cette
nature, d'où se déduirait l'extension vectorielle universelle par
« enveloppe vectorielle ».\footnote{La phrase de la page 210 est très
incertaine ; seules les formules sont garanties.}

\pagerange{214}{217}
\section*{XVII. $c^1(L)$ et connexions (pages 213 à 217)}

Le titre est celui de la chemise (page 213), lecture incertaine du mot
« connexions ». Pour $X/S$ lisse, soit $\Omega^{\geq 1}_{X/S}$ le complexe de
De~Rham tronqué. La classe de $L$ dans $\mathbb{H}^2(X, \Omega^{\geq 1}_{X/S})$,
image de $L$ par $d\log : \mathcal{O}^* \to \Omega^{\geq 1}[1]$, est
l'obstruction à l'existence d'une connexion intégrable sur $L$ ; par la
suite spectrale elle se décompose en ⓐ la classe de Hodge
$c^1_{\mathrm{Hdg}}(L) \in H^1(X, \Omega^1)$, et ⓑ une seconde obstruction dans
$Z^2/\mathrm{Im}$. En caractéristique $p$, l'obstruction est tuée par $p$ :
elle s'annule sur $p\,\mathrm{Pic}(X)$ et, « à vérifier », sur
$\mathrm{Pic}^\tau$ si $X$ est propre.

Quand elle est nulle, les connexions intégrables diffèrent par des formes
fermées $\omega$, et la $p$-courbure $v$ — le « défaut de compatibilité aux
puissances $p$-ièmes », section de $\Omega^{1\,(p)}$ — devient
\[
  v' = v + \bigl(\omega^{\otimes p} - F^*_{X/S}\widetilde C\omega\bigr) ,
\]
d'où une nouvelle obstruction dans $\Gamma(\Omega^{1\otimes p})/\mathrm{Im}
\bigl(Z^1 \to \Gamma(\Omega^{1\otimes p})\bigr)$, l'application $\omega \mapsto
\omega^{\otimes p} - F^*\widetilde C\omega$ n'étant « pas linéaire ! ». Sous des
hypothèses de platitude (« OK si $X$ de dimension relative 1 »), le foncteur
des modules inversibles à connexion intégrable serait représentable par un
fibré vectoriel au-dessus de $\underline{\mathrm{Pic}}$ — la page s'arrête sur
cette phrase. Les pages 215 et 217, barrées, portent sur des catégories de
couples $S \leftarrow U_0 \hookrightarrow U$ (voisinages infinitésimaux) et leur
ind-représentabilité ; elles sont à peine lisibles.

\pagerange{219}{220}
\section*{XVIII. Connexion canonique de Gauss--Manin : un plan (pages 218 à 220)}

La chemise (page 218) porte « Connexion canonique de Gauss-Manin ». Un plan :
1) compatibilités pour $S$, $T$, $E$ variables ; 2) complexes de De~Rham
tronqués, filtration à deux crans, suites exactes de Hodge--De~Rham, et
« transversalités de Griffiths et de Deligne » ; 3) variantes stratifiées en
caractéristique nulle, et variantes cristallines en caractéristique
quelconque. La page 220 commence une construction : pour un groupe $G$ de
faisceau de sections $\underline\Gamma(G)$ et d'algèbre de Lie $\mathfrak{g}$,
l'application $g \mapsto g^{-1}dg$ (forme de Maurer--Cartan) vers les
$1$-formes fermées à coefficients dans $\mathfrak{g}$, composée avec la
factorisation $\tilde D$ d'un opérateur différentiel $D : \mathfrak{g} \to
\mathcal{M}$ par $d^0$ (« spécial à car.~0 ? ») ; la page s'arrête là.

\pagerange{221}{229}
\section*{XIX. Intégrales de formes fermées, et équations de Picard--Fuchs (pages 221 à 229)}

Soit $f : X \to S$ une submersion holomorphe, $\omega$ une $1$-forme
relative fermée, $a$, $b$ deux sections holomorphes avec $a(s)$, $b(s)$ dans
la même composante de $X_s$. Localement sur $S$, des arcs $C_s$ de $a(s)$ à
$b(s)$ dépendant continûment de $s$ définissent une fonction holomorphe
$I(s) = \int_{a(s)}^{b(s)} \omega_s$, déterminée modulo les périodes
$\int_{Z_s}\omega_s$ sur les $1$-cycles. Une famille $\mu = (D_i, \omega_i)$
d'opérateurs différentiels $D_i$ sur $S$ et de formes fermées \emph{définit
une équation de Picard--Fuchs} si $\sum_i D_i\int_{Z_s}\omega_{i,s} = 0$ pour
toute famille continue de cycles ; alors
\[
  \mu(a,b) = \sum_i D_{i,s} \int_{a(s)}^{b(s)} \omega_{i,s}
\]
est une fonction holomorphe bien définie, et $\mu(a,b) + \mu(b,c) = \mu(a,c)$.

\textbf{Critère.} Soit $\mathcal{F} = f_*(Z\Omega^1_{X/S})/d f_*\mathcal{O}_X$, muni de
sa connexion intégrable naturelle (par la formule de Cartan ; « attention,
en réalité il faut que $X \to S$ soit steinien »). Il suffit que
$\sum_i D_i\,\omega'_i = 0$ dans $\mathcal{F}$, $\omega'_i$ la classe de
$\omega_i$. En effet, si $X/S$ porte localement une connexion intégrable
(un feuilletage transverse aux fibres), on relève les $D_i$ en opérateurs
$\bar D_i$ commutant à $d_{X/S}$, et la condition dit que $\varpi = \sum \bar
D_i\omega_i = d\varphi$ ; la formule de dérivation sous le signe somme
\[
  D_s\int_{a(s)}^{b(s)}\omega_s = \int_{a(s)}^{b(s)} (\bar D\omega)_s ,
\]
valable quand $a$, $b$ sont des sections \emph{horizontales}, et qui se
vérifie sur de petits chemins où $\omega = d\varphi$, donne alors
$\sum_i D_i \int_{Z_s}\omega_i = \int_{Z_s}d\varphi = 0$. Réciproquement,
admettant la formule $D_s\int_{Z_s}\omega' = \int_{Z_s}(\bar D\omega')_s$
(15), le théorème de De~Rham sur les fibres montre que la condition est
nécessaire.\footnote{La formule (15) est dite « la propriété transcendante
essentielle de la connexion intégrable canonique sur $\mathcal{F}$ » ; elle est
admise, non démontrée. La lettre cursive que la transcription note
$\mathcal{F}$ est de forme incertaine.}

\textbf{Calcul algébrique de $\mu(a,b)$.} Localement, si $a$ et $b$ sont
horizontales, $\mu(a,b) = \varphi(b(s)) - \varphi(a(s))$, $\varphi$ défini
par $\sum\bar D_i\omega_i = d\varphi$ ; en général, avec les sections
horizontales $a_0$, $b_0$ qui coïncident avec $a$, $b$ en $s_0$,
$\mu(a,b) = \mu(a_0,b_0) + \mu(b_0,b) - \mu(a_0,a)$, et il reste un calcul
purement local, que la page commence puis barre. C'est l'avatar analytique
de la « section de Manin » de la station XI.

\pagerange{230}{232}
\section*{XX. Opérateurs différentiels à coefficients constants (pages 230 à 232)}

Pour $S \to T$ : 1) si $F \xrightarrow{D} G \xrightarrow{D'} H$ est une suite
d'opérateurs différentiels d'ordre $1$ « formellement exacte », les suites
$\mathrm{Diff}^i(P,F) \to \mathrm{Diff}^{i+1}(P,G) \to \mathrm{Diff}^{i+2}(P,H)$
(composition à gauche) et $\mathrm{Diff}^i(H,P) \to \mathrm{Diff}^{i+1}(G,P) \to
\mathrm{Diff}^{i+2}(F,P)$ (à droite) sont exactes. 2) Si $E$ est localement libre
stratifié et que le lemme de Poincaré formel vaut pour $E \otimes
\Omega^\bullet_{S/T}$ (par exemple $S$ lisse sur $T$ de caractéristique nulle),
la suite
\[
  0 \leftarrow \mathrm{Hom}(E,P) \leftarrow \mathrm{Diff}^i(E,P)
  \xleftarrow{\,\circ d^0_E\,} \mathrm{Diff}^{i-1}(E \otimes \Omega^1, P)
  \xleftarrow{\,\circ d^1_E\,} \cdots
\]
est exacte. Le noyau $\mathrm{Diff}^i_+(E,P)$ de la première flèche — les
opérateurs « à coefficients constants » relativement à la stratification —
est donc $\mathrm{Diff}^{i-1}(E\otimes\Omega^1,P)/\mathrm{Im}\,\mathrm{Diff}^{i-2}(E
\otimes\Omega^2,P)$, ce qui précise le lemme de la page 159, et il s'envoie
(injectivement ?) dans les opérateurs différentiels de $\operatorname{Ker} d^1_E$
dans $P$ ; le module universel est $P^i(E)/\mathcal{O}_S\cdot\delta_i(E)$. 3) Pour
un sous-module $N \subset E$ et $M = E/N$, les opérateurs sur $M$ qui
proviennent d'opérateurs à coefficients constants sur $E$ sont représentés
par $P^i(M)/\mathcal{O}_S\cdot\rho_i(E)$, $\rho_i : E \to P^i(E) \to P^i(M)$ ; ils
sont déterminés par $\rho_\infty : E \to P^\infty(M)$, soumis à des conditions
d'intégrabilité. C'est le cadre des opérateurs de Picard--Fuchs de la
station XI, avec $M = E''$.

\pagerange{233}{235}
\section*{XXI. Construction de la connexion canonique en cohomologie de De~Rham (pages 233 à 235)}

Les pages sont numérotées par lui « 1 » et « 2 », sous ce titre encadré.

Soit $T_0 \hookrightarrow T$ une immersion de carré nul, d'idéal $I$, $X_0$ lisse
sur $T_0$, et $X'$, $X''$ deux relèvements lisses de $X_0$ sur $T$
(localement isomorphes). Le faisceau $\mathcal{P}$ des germes de
$T$-isomorphismes $X' \to X''$ induisant l'identité sur $X_0$ est un torseur
sous le faisceau abélien $\mathcal{G} = \underline{\mathrm{Hom}}(\Omega^1_{X_0/T_0},
f_0^*I) = T_{X_0/T_0} \otimes f_0^*I$. Chaque $p \in \mathcal{P}$ induit un
isomorphisme de complexes $u(p)$ entre $\Omega^\bullet_{X'/T}$ et
$\Omega^\bullet_{X''/T}$, et chaque $g \in \mathcal{G}$, champ de vecteurs à
valeurs dans $I$, donne par produit intérieur un opérateur $k(g)$ de degré
$-1$, avec la formule d'homotopie de Cartan
\[
  u(p + g) - u(p) = d\,k(g) + k(g)\,d .
\]
\footnote{La page 233 écrit $u(pg) = u(p)v(g)$ et $v(g) - \mathrm{id} = dk(g) +
k(g)d$, la page 234 $v(g) = dk(g) + k(g)d$ ; $\mathcal{G}$ étant abélien, on
écrit son action additivement et l'on combine les deux relations. Le sens
des flèches $u(p)$ (image directe des formes par $p$, ou image inverse par
$p^{-1}$) n'est pas fixé par la page ; il est sans effet sur ce qui suit.}

\textbf{La construction générale} (pages 234 et 235). Sur un espace annelé,
soient $K'$, $K''$ deux complexes, $\mathcal{G}$ un faisceau abélien,
$\mathcal{P}$ un $\mathcal{G}$-torseur, $u : \mathcal{P} \to
\underline{\mathrm{Hom}}^0(K', K'')$ et $k : \mathcal{G} \to
\underline{\mathrm{Hom}}^{-1}(K',K'')$ avec $u(p+g) - u(p) = dk(g) + k(g)d$. Soit
$\overline{\mathcal{G}}$ l'extension $0 \to \mathcal{G} \to \overline{\mathcal{G}}
\xrightarrow{j} \mathbb{Z} \to 0$ telle que $j^{-1}(1) = \mathcal{P}$, et
$\mathcal{L}$ le complexe $\mathcal{G} \to \overline{\mathcal{G}}$ (en degrés $-1$
et $0$), quasi-isomorphe à $\mathbb{Z}$. Les données $k$ et $u$ définissent un
morphisme de complexes $\varphi : \mathcal{L} \to \mathcal{H} =
\underline{\mathrm{Hom}}^\bullet(K',K'')$, avec $\varphi^{-1} = k$ et
$\varphi^0|_{\mathcal{P}} = u$ — un morphisme de $\overline{\mathcal{G}}$ dans un
faisceau étant la donnée d'un morphisme sur $\mathcal{G}$ et d'une application
sur $\mathcal{P}$ compatibles. D'où, dans la catégorie dérivée,
\[
  \mathbb{Z} \simeq \mathcal{L} \longrightarrow \mathcal{H} \longrightarrow
  \mathbf{R}\underline{\mathrm{Hom}}(K',K''),
\]
puis $\mathbf{R}\Gamma_X(\mathcal{O}_X) \to \mathbf{R}\underline{\mathrm{Hom}}(K',K'')$, et
en degré $0$ un élément de $\mathrm{Hom}_{D}(K', K'')$ : un morphisme canonique
$K' \to K''$ dans la catégorie dérivée, sans choix de $p$.\footnote{Vérifié
en Lean (lean/Grothendieck/Folder93.lean, \texttt{j\_eq\_zero\_iff},
\texttt{j\_eq\_one\_iff}, \texttt{$\delta$\_$\varphi$0},
\texttt{quotient\_map\_eq}, \texttt{derived\_map\_eq}) dans le cas d'un
torseur sous un groupe abélien $G$, $K'$ et $K''$ étant des complexes d'une
catégorie préadditive quelconque. L'extension $\overline{G}$ se construit sans
choisir de point, comme quotient de $G \times \mathbb{Z}^{(\mathcal{P})}$ par les
relations $[p + g] - [p] = g$ ; $[G \to \overline{G}]$ est quasi-isomorphe à
$\mathbb{Z}$ ; $\varphi^0$, égal à $u$ sur $\mathcal{P}$ et à $dk + kd$ sur $G$,
fait de $(k, \varphi^0)$ un morphisme de complexes ; et la classe de $u(p)$ dans
la catégorie d'homotopie, et dans la catégorie dérivée, ne dépend pas de $p$,
l'homotopie entre $u(p')$ et $u(p)$ étant $k(p' - p)$. Pour un groupe, le
torseur a toujours un point et l'énoncé se réduit à cette indépendance ; le cas
visé par la page, un torseur sous un faisceau sans section globale, n'est pas
vérifié, ni l'application aux complexes de De~Rham.}

Appliqué à $K' = \Omega^\bullet_{X'/T}$, $K'' = \Omega^\bullet_{X''/T}$, c'est
l'isomorphisme canonique entre les cohomologies de De~Rham de deux
relèvements — le fait que la cohomologie de De~Rham est un \emph{cristal} —,
dont la connexion de Gauss--Manin se déduit en l'appliquant aux deux
relèvements $\mathrm{pr}_1^*$, $\mathrm{pr}_2^*$ au voisinage infinitésimal de
la diagonale. Cette dernière étape n'est pas dans le dossier : la page 235
s'arrête sur « un homom.\ de la cat.\ dérivée ».\footnote{La déduction de la
connexion est notre lecture du titre, « Construction de la connexion
canonique ». Le procédé est celui de son exposé « Crystals and the de Rham
cohomology of schemes » (1968) et de la thèse de P.~Berthelot.}

\end{document}
